%======================================================================
% Chapter 10: Statistics for Machine Learning
%======================================================================
\chapter{Statistics for Machine Learning}\label{ch:statistics}

\begin{motivation}
Probability starts from a known distribution and predicts what data it will
produce. Statistics runs in the opposite direction: it starts from data and
draws conclusions about the unknown distribution that produced them. That is
exactly the situation of machine learning. We never know the true relationship
between inputs and outputs; we only have a finite, noisy sample of examples,
and from it we must
\begin{itemize}
  \item \emph{estimate} the parameters of a model (training);
  \item judge how reliable a number computed from a finite sample is (is a test
    accuracy of $91.2\%$ really better than $90.8\%$?);
  \item understand why the standard loss functions, squared error and
    cross-entropy, are what they are;
  \item understand why a model that fits its training data perfectly can fail
    badly on new data, and what to do about it.
\end{itemize}
This chapter covers these ideas: estimators and their errors, the law of large
numbers and the central limit theorem, maximum likelihood and Bayesian
estimation, hypothesis testing, and generalisation.
\end{motivation}

\begin{objectives}
\begin{itemize}
  \item summarise data with means, medians, variances and percentiles;
  \item describe estimators by their bias, variance and standard error;
  \item state and use the law of large numbers and Monte Carlo estimation;
  \item state the central limit theorem and compute confidence intervals;
  \item derive maximum likelihood estimates;
  \item derive the squared-error and cross-entropy losses from likelihoods;
  \item connect priors with regularisation through MAP estimation;
  \item carry out and correctly interpret a hypothesis test;
  \item explain overfitting, underfitting and the bias--variance trade-off, and
    use validation sets and cross-validation.
\end{itemize}
\end{objectives}

%----------------------------------------------------------------------
\section{Data, samples and summary statistics}\label{sec:descriptive}
%----------------------------------------------------------------------

Statistics views a data set $x_1, \dots, x_n$ as the observed values of random
variables $X_1, \dots, X_n$. Usually they are assumed to be \emph{i.i.d.}
(\cref{sec:independence-events}): drawn independently from one and the same
distribution, called the \term{population}. The data are a \term{sample} from
the population.

Numbers that describe the population, such as its mean $\mu$ and variance
$\sigma^2$, are called \term[parameter (statistics)]{parameters}; they are fixed
but unknown. Numbers computed from the sample are called
\term[statistic]{statistics}; they are known, but they would come out
differently for a different sample.

\begin{definition}[Sample mean, variance and standard deviation]\label{def:sample-stats}
For data $x_1, \dots, x_n$, the \term{sample mean} and the \term{sample variance}
are
\[
  \bar x = \frac1n\sum_{i=1}^n x_i,
  \qquad
  s^2 = \frac{1}{n-1}\sum_{i=1}^n (x_i - \bar x)^2 ,
\]
and the \emph{sample standard deviation} is $s = \sqrt{s^2}$.
\end{definition}

The division by $n - 1$ instead of $n$ looks strange; \cref{sec:estimators}
explains it. For large $n$ the difference is negligible.

\begin{example}[Computing summary statistics]\label{ex:summary-stats}
Find the sample mean, variance and standard deviation of the data
$2, 4, 4, 4, 5, 5, 7, 9$.

\sol The mean is $\bar x = \frac{40}{8} = 5$. The deviations from the mean are
$-3$, $-1$, $-1$, $-1$, $0$, $0$, $2$ and $4$, and their squares add up to
$9 + 1 + 1 + 1 + 0 + 0 + 4 + 16 = 32$. So $s^2 = \frac{32}{7} \approx 4.57$ and
$s \approx 2.14$.
\end{example}

\subsection{Medians and percentiles}

\begin{definition}[Median and percentiles]\label{def:median}
The \term{median} of a data set is its middle value after sorting (the average
of the two middle values if $n$ is even). More generally, the
$p$-th \term{percentile} is a value below which $p\%$ of the data lie; the
median is the $50$th percentile, and the $25$th and $75$th percentiles are
called the lower and upper \term[quartile]{quartiles}.
\end{definition}

\begin{example}[Mean versus median]\label{ex:mean-median}
The annual salaries (in thousands) of eight employees are
\[
  32,\ 35,\ 38,\ 40,\ 41,\ 44,\ 47,\ 250 .
\]
The mean is $\frac{527}{8} \approx 65.9$, more than all but one of the salaries.
The median is $\frac{40 + 41}{2} = 40.5$, a much better description of a typical
salary.
\end{example}

\begin{intuition}[Two kinds of centre]
The mean is the balance point of the data, so one extreme value (an
\term{outlier}) can drag it far away. The median only cares about the order of
the values and is \emph{robust} to outliers. There is a deeper connection: the
mean is the number $c$ that minimises the total \emph{squared} distance
$\sum_i(x_i - c)^2$ (\cref{ex:mean-minimises}), while the median minimises the
total \emph{absolute} distance $\sum_i\abs{x_i - c}$ (an exercise below). This is
the same contrast as between the squared and the absolute error
(\cref{sec:absolute-value}).
\end{intuition}

\begin{mlnote}[Look at your data]
Before training anything, practitioners compute summary statistics of every
feature: means and standard deviations (for standardisation,
\cref{sec:variance}), minima and maxima (to find impossible values), and
percentiles (to see skew and outliers). Percentiles are also how systems are
measured: the ``p99 latency'' of a model server is the $99$th percentile of its
response times.
\end{mlnote}

\sectionexercises

\begin{exercise}
For the data $4, 8, 6, 5, 3, 10$, compute the mean, the median, the sample
variance and the sample standard deviation.
\end{exercise}
\begin{solution}
Mean $\tfrac{36}{6} = 6$. Sorted: $3, 4, 5, 6, 8, 10$, so the median is
$\tfrac{5 + 6}{2} = 5.5$. Deviations $-2, 2, 0, -1, -3, 4$; squares add up to
$34$, so $s^2 = \tfrac{34}{5} = 6.8$ and $s \approx 2.61$.
\end{solution}

\begin{exercise}
In the previous exercise, the value $10$ is replaced by $100$ (a typo). What are
the new mean and median?
\end{exercise}
\begin{solution}
Mean $\tfrac{126}{6} = 21$; the median is unchanged at $5.5$.
\end{solution}

\begin{exercise}
Every value of a data set is transformed by $y_i = 2x_i + 3$. How do the mean,
the median and the standard deviation change?
\end{exercise}
\begin{solution}
$\bar y = 2\bar x + 3$ and the median becomes $2m + 3$ (the order of the values is
unchanged). The deviations from the mean are doubled and the shift cancels, so
$s_y = 2s_x$.
\end{solution}

\begin{exercise}
Let $g(c) = \sum_{i=1}^n\abs{x_i - c}$. Explain why, if $c$ is not a data value,
the slope of $g$ at $c$ equals (number of data values below $c$) minus (number
of data values above $c$). Deduce that $g$ is smallest at the median.
\end{exercise}
\begin{solution}
Each term $\abs{x_i - c}$ has slope $+1$ in $c$ if $x_i < c$ (moving $c$ right
moves it away from $x_i$) and $-1$ if $x_i > c$. Adding up gives the stated
slope. It is negative while fewer than half of the values lie below $c$ and
positive once more than half do, so $g$ decreases up to the median and increases
after it.
\end{solution}

%----------------------------------------------------------------------
\section{Estimators: bias, variance and standard error}\label{sec:estimators}
%----------------------------------------------------------------------

\begin{definition}[Estimator]\label{def:estimator}
An \term{estimator} is a rule that computes, from the data, a guess $\hat\theta$
(read ``theta hat'') of an unknown parameter $\theta$. Because the data are
random, an estimator is a random variable, and its distribution is called its
\term{sampling distribution}.
\end{definition}

The sample mean $\bar X$ is an estimator of $\mu$; a model's accuracy on a test
set is an estimator of its accuracy on all possible inputs; the weights found by
training are estimators of the ``best'' weights. How good is an estimator? Two
numbers answer this.

\begin{definition}[Bias and unbiasedness]\label{def:bias}
The \term{bias} of an estimator $\hat\theta$ is
$\Bias(\hat\theta) = \E[\hat\theta] - \theta$. The estimator is
\term[unbiased estimator]{unbiased} if its bias is $0$: on average, it hits the
true value.
\end{definition}

\begin{theorem}[Mean and variance of the sample mean]\label{thm:sample-mean}
Let $X_1, \dots, X_n$ be i.i.d.\ with mean $\mu$ and variance $\sigma^2$. Then
\[
  \E[\bar X] = \mu
  \qquad\text{and}\qquad
  \Var(\bar X) = \frac{\sigma^2}{n} .
\]
\end{theorem}

\begin{proof}
By linearity of expectation, $\E[\bar X] = \frac1n\sum_i\E[X_i] = \frac1n\cdot
n\mu = \mu$. By independence the variances of the $X_i$ add up, so
$\Var\left(\sum_iX_i\right) = n\sigma^2$, and by \cref{thm:variance-properties}
$\Var(\bar X) = \frac{1}{n^2}\,n\sigma^2 = \frac{\sigma^2}{n}$.
\end{proof}

\begin{definition}[Standard error]\label{def:standard-error}
The \term{standard error} of an estimator is its standard deviation. For the
sample mean it is $\operatorname{SE}(\bar X) = \sigma/\sqrt n$, which we
estimate by $s/\sqrt n$.
\end{definition}

\begin{intuition}[The square-root law]
Averaging $n$ independent measurements makes the error $\sqrt n$ times smaller,
not $n$ times. To halve the uncertainty you need four times as much data; to
gain one more correct decimal digit, a hundred times as much. This law explains
a great deal about machine learning, from the noise in mini-batch gradients to
the size of test sets.
\end{intuition}

\begin{example}[How precise is a test accuracy?]\label{ex:accuracy-se}
A classifier is correct on $850$ of $1000$ test examples. How precise is the
accuracy estimate $0.85$?

\sol Each test example gives a Bernoulli variable, $1$ if the prediction is
correct and $0$ otherwise, with some unknown mean $p$ (the true accuracy) and
variance $p(1 - p)$. The test accuracy is their sample mean, so its standard
error is
\[
  \sqrt{\frac{p(1 - p)}{n}} \approx \sqrt{\frac{0.85\times0.15}{1000}}
  \approx 0.011 .
\]
The accuracy is uncertain by about one percentage point, so a difference
between $85.3\%$ and $85.7\%$ on this test set means very little.
\end{example}

\subsection{Why divide by \texorpdfstring{$n - 1$}{n-1}?}

\begin{theorem}[An unbiased estimator of the variance]\label{thm:bessel}
For i.i.d.\ $X_1, \dots, X_n$ with variance $\sigma^2$,
\[
  \E\left[\sum_{i=1}^n(X_i - \bar X)^2\right] = (n - 1)\,\sigma^2 .
\]
So dividing by $n$ gives an estimator that is too small on average, by the
factor $\frac{n-1}{n}$, whereas $s^2$, which divides by $n - 1$, is unbiased.
\end{theorem}

\begin{proof}
Write $X_i - \bar X = (X_i - \mu) - (\bar X - \mu)$ and expand. Since
$\sum_i(X_i - \mu) = n(\bar X - \mu)$,
\[
  \sum_i(X_i - \bar X)^2 = \sum_i(X_i - \mu)^2 - 2n(\bar X - \mu)^2 + n(\bar X - \mu)^2
  = \sum_i(X_i - \mu)^2 - n(\bar X - \mu)^2 .
\]
Taking expectations, using $\E[(X_i - \mu)^2] = \sigma^2$ and
$\E[(\bar X - \mu)^2] = \Var(\bar X) = \sigma^2/n$, gives
$n\sigma^2 - \sigma^2 = (n - 1)\sigma^2$.
\end{proof}

The intuition: the data are on average closer to their own mean $\bar x$ (which
was computed from them) than to the true mean $\mu$, so squared deviations from
$\bar x$ underestimate the spread. Dividing by $n - 1$, a correction known as
\term{Bessel's correction}, compensates exactly.

\subsection{Mean squared error}

Bias is not everything: an unbiased estimator that is wildly variable is useless.
The overall quality of an estimator combines both.

\begin{theorem}[Bias--variance decomposition of the MSE]\label{thm:mse-decomposition}
The \term[mean squared error!of an estimator]{mean squared error} of an
estimator, $\MSE(\hat\theta) = \E\bigl[(\hat\theta - \theta)^2\bigr]$, satisfies
\[
  \keyeq{\MSE(\hat\theta) = \Bias(\hat\theta)^2 + \Var(\hat\theta)}
\]
\end{theorem}

\begin{proof}
Let $m = \E[\hat\theta]$. Then $\hat\theta - \theta = (\hat\theta - m) + (m -
\theta)$, where the second part is a constant. Squaring and taking expectations,
\[
  \E\bigl[(\hat\theta - \theta)^2\bigr] = \E\bigl[(\hat\theta - m)^2\bigr]
  + 2(m - \theta)\,\E[\hat\theta - m] + (m - \theta)^2 .
\]
The middle term vanishes because $\E[\hat\theta - m] = 0$; the other two are the
variance and the squared bias.
\end{proof}

\begin{mlnote}[A little bias can help]
A biased estimator can have a smaller mean squared error than an unbiased one, if
it has much smaller variance. Shrinking an estimate towards zero introduces
bias but reduces variance, and this is exactly what regularisation does to the
weights of a model (\cref{ex:ridge-1d}). We meet the same trade-off for whole
models in \cref{sec:generalisation}.
\end{mlnote}

\sectionexercises

\begin{exercise}
Measurements have standard deviation $\sigma = 10$. How many must be averaged for
the standard error of the mean to be $1$? To be $0.5$?
\end{exercise}
\begin{solution}
$\tfrac{10}{\sqrt n} = 1$ gives $n = 100$; $\tfrac{10}{\sqrt n} = 0.5$ gives
$n = 400$.
\end{solution}

\begin{exercise}
From i.i.d.\ data $X_1, \dots, X_{25}$ with mean $\mu$ and variance $\sigma^2$,
someone estimates $\mu$ by the first observation $X_1$ alone. Is this estimator
unbiased? Compare its variance with that of $\bar X$.
\end{exercise}
\begin{solution}
$\E[X_1] = \mu$, so it is unbiased, but its variance is $\sigma^2$, versus
$\tfrac{\sigma^2}{25}$ for $\bar X$. Being unbiased is not enough; $\bar X$ uses
all the data and is much more precise.
\end{solution}

\begin{exercise}
With $n = 5$ observations, what is the expected value of the variance estimator
$\frac1n\sum_i(X_i - \bar X)^2$, and what is its bias?
\end{exercise}
\begin{solution}
By \cref{thm:bessel} its expectation is $\tfrac{n-1}{n}\sigma^2 = 0.8\sigma^2$, so
its bias is $-0.2\sigma^2$.
\end{solution}

\begin{exercise}
Estimator A is unbiased with variance $4$. Estimator B has bias $1$ and variance
$1$. Which has the smaller mean squared error?
\end{exercise}
\begin{solution}
$\MSE(A) = 0 + 4 = 4$ and $\MSE(B) = 1 + 1 = 2$, so the biased estimator B is
better on average.
\end{solution}

%----------------------------------------------------------------------
\section{The law of large numbers and Monte Carlo}\label{sec:lln}
%----------------------------------------------------------------------

Why should averages of data tell us anything about the population? Because of
the following theorem, which turns the intuition ``averages settle down'' into
mathematics.

\begin{theorem}[Law of large numbers]\label{thm:lln}
Let $X_1, X_2, \dots$ be i.i.d.\ with mean $\mu$ and finite variance $\sigma^2$,
and let $\bar X_n = \frac1n\sum_{i=1}^nX_i$. Then for every $\varepsilon > 0$,
\[
  P\bigl(\abs{\bar X_n - \mu} \ge \varepsilon\bigr) \to 0
  \qquad\text{as } n \to \infty .
\]
\end{theorem}
\index{law of large numbers}

\begin{proof}
Apply Chebyshev's inequality (\cref{sec:variance}) to $\bar X_n$, whose mean is
$\mu$ and variance $\frac{\sigma^2}{n}$ (\cref{thm:sample-mean}):
$P\bigl(\abs{\bar X_n - \mu} \ge \varepsilon\bigr) \le \frac{\sigma^2}{n\varepsilon^2}$,
which tends to $0$.
\end{proof}

\begin{figure}[htbp]
\centering
\begin{tikzpicture}
\begin{axis}[mlplot,axis lines=left,xmode=log,log basis x=10,width=11cm,height=6cm,
  xmin=1,xmax=10000,ymin=1,ymax=6.3,ytick={1,...,6},clip=false,
  xtick={1,10,100,1000,10000},xticklabels={$1$,$10$,$100$,$1000$,$10\,000$},
  xlabel={number of rolls $n$ (logarithmic scale)},ylabel={running average},
  xlabel style={at={(ticklabel cs:0.5)},anchor=north},
  ylabel style={at={(ticklabel cs:0.5)},anchor=south,rotate=90}]
  \draw[gray,dashed,thick] (axis cs:1,3.5) -- (axis cs:10000,3.5);
  \node[gray,font=\small,anchor=west] at (axis cs:10000,3.5) {mean $3.5$};
  \addplot[PlotBlue] coordinates {(1,6.000) (2,6.000) (3,5.667) (4,5.750) (5,4.800) (6,4.833) (7,5.000) (8,4.625) (9,4.333)
    (10,4.000) (12,4.250) (13,4.231) (15,4.267) (16,4.062) (18,4.278) (21,4.286) (23,4.130)
    (26,4.115) (29,4.034) (33,4.121) (37,4.135) (42,4.024) (47,4.000) (53,4.019) (59,4.017)
    (66,3.924) (75,3.853) (84,3.714) (94,3.755) (106,3.849) (119,3.832) (134,3.761)
    (150,3.687) (169,3.663) (190,3.726) (213,3.681) (240,3.708) (269,3.732) (303,3.647)
    (340,3.638) (382,3.644) (429,3.641) (483,3.588) (542,3.594) (609,3.573) (685,3.531)
    (769,3.506) (864,3.497) (971,3.454) (1091,3.461) (1226,3.445) (1378,3.459) (1548,3.467)
    (1740,3.464) (1955,3.471) (2197,3.471) (2468,3.490) (2774,3.480) (3117,3.479) (3502,3.470)
    (3935,3.472) (4421,3.487) (4968,3.487) (5583,3.486) (6273,3.499) (7049,3.506) (7920,3.511)
    (8900,3.502) (10000,3.502)};
  \addplot[PlotRed] coordinates {(1,3.000) (2,3.000) (3,3.667) (4,3.000) (5,3.200) (6,2.833) (7,2.571) (8,2.750) (9,3.111)
    (10,3.400) (12,3.500) (13,3.385) (15,3.467) (16,3.438) (18,3.500) (21,3.524) (23,3.652)
    (26,3.462) (29,3.379) (33,3.485) (37,3.378) (42,3.452) (47,3.404) (53,3.302) (59,3.203)
    (66,3.212) (75,3.160) (84,3.131) (94,3.170) (106,3.255) (119,3.218) (134,3.336)
    (150,3.313) (169,3.361) (190,3.411) (213,3.418) (240,3.375) (269,3.413) (303,3.403)
    (340,3.444) (382,3.421) (429,3.422) (483,3.439) (542,3.472) (609,3.471) (685,3.477)
    (769,3.471) (864,3.454) (971,3.458) (1091,3.459) (1226,3.495) (1378,3.498) (1548,3.477)
    (1740,3.494) (1955,3.495) (2197,3.512) (2468,3.500) (2774,3.494) (3117,3.478) (3502,3.482)
    (3935,3.482) (4421,3.467) (4968,3.482) (5583,3.484) (6273,3.484) (7049,3.479) (7920,3.480)
    (8900,3.483) (10000,3.490)};
  \addplot[PlotGreen] coordinates {(1,1.000) (2,1.000) (3,1.000) (4,1.750) (5,2.000) (6,2.667) (7,2.571) (8,2.750) (9,2.889)
    (10,2.700) (12,2.917) (13,3.077) (15,3.200) (16,3.125) (18,3.167) (21,3.190) (23,3.174)
    (26,3.308) (29,3.172) (33,3.152) (37,3.324) (42,3.214) (47,3.191) (53,3.283) (59,3.322)
    (66,3.303) (75,3.200) (84,3.143) (94,3.202) (106,3.377) (119,3.345) (134,3.366)
    (150,3.313) (169,3.325) (190,3.316) (213,3.291) (240,3.300) (269,3.297) (303,3.314)
    (340,3.350) (382,3.377) (429,3.380) (483,3.402) (542,3.408) (609,3.374) (685,3.412)
    (769,3.404) (864,3.418) (971,3.434) (1091,3.475) (1226,3.481) (1378,3.499) (1548,3.490)
    (1740,3.495) (1955,3.490) (2197,3.508) (2468,3.522) (2774,3.514) (3117,3.515) (3502,3.505)
    (3935,3.500) (4421,3.502) (4968,3.486) (5583,3.485) (6273,3.491) (7049,3.494) (7920,3.494)
    (8900,3.496) (10000,3.502)};
\end{axis}
\end{tikzpicture}
\caption{The law of large numbers. Each curve shows the average of the first $n$
  rolls of a die, for three independent sequences of $10\,000$ rolls. The
  averages wander at first and then settle down to the mean $3.5$.}
\label{fig:lln}
\end{figure}

\begin{pitfall}[The gambler's fallacy]
The law of large numbers does \emph{not} say that after a run of heads, tails
become ``due''. A coin has no memory: after ten heads in a row, the next toss is
still heads with probability $\frac12$. Early imbalances are not compensated;
they are \emph{diluted} by the many tosses that follow.
\end{pitfall}

\subsection{Monte Carlo estimation}

The law of large numbers gives a universal recipe for computing expectations,
probabilities and integrals: \emph{simulate and average}.

\begin{definition}[Monte Carlo estimation]\label{def:monte-carlo}
To estimate $\E[g(X)]$, draw independent samples $x_1, \dots, x_n$ from the
distribution of $X$ and compute
\[
  \E[g(X)] \approx \frac1n\sum_{i=1}^n g(x_i) .
\]
This is the \term{Monte Carlo method}. Its standard error is
$\sigma_g/\sqrt n$, where $\sigma_g$ is the standard deviation of $g(X)$,
\emph{whatever the dimension of $X$}.
\end{definition}

\begin{example}[Estimating $\pi$ by throwing darts]\label{ex:monte-carlo-pi}
Choose points uniformly at random in the unit square $[0, 1]^2$. The probability
that a point lands inside the quarter disc $x^2 + y^2 \le 1$ is its area,
$\frac\pi4$. In a simulation with $400$ points (\cref{fig:monte-carlo-pi}), $318$
landed inside, giving the estimate
\[
  \pi \approx 4\times\frac{318}{400} = 3.18 .
\]
The standard error is about $4\sqrt{0.785\times0.215/400} \approx 0.08$,
and the estimate is indeed within $0.04$ of $\pi$.
\end{example}

\begin{figure}[htbp]
\centering
\begin{tikzpicture}
\begin{axis}[mlplot,axis lines=box,width=6.2cm,height=6.2cm,xmin=0,xmax=1,ymin=0,ymax=1,
  axis equal image,xtick={0,0.5,1},ytick={0,0.5,1},grid=none]
  \addplot[only marks,mark=*,mark size=0.7pt,PlotBlue] coordinates {(0.262,0.298) (0.814,0.092) (0.600,0.729) (0.188,0.055) (0.275,0.657) (0.562,0.150)
    (0.433,0.669) (0.423,0.633) (0.392,0.187) (0.346,0.511) (0.318,0.924) (0.471,0.694)
    (0.107,0.105) (0.202,0.884) (0.644,0.407) (0.517,0.593) (0.862,0.438) (0.829,0.498)
    (0.693,0.339) (0.523,0.216) (0.101,0.039) (0.702,0.456) (0.592,0.766) (0.407,0.196)
    (0.172,0.181) (0.604,0.113) (0.020,0.833) (0.099,0.451) (0.488,0.620) (0.750,0.574)
    (0.617,0.507) (0.965,0.227) (0.689,0.555) (0.042,0.296) (0.013,0.297) (0.010,0.827)
    (0.110,0.057) (0.318,0.049) (0.390,0.366) (0.523,0.007) (0.148,0.210) (0.441,0.302)
    (0.613,0.285) (0.059,0.209) (0.563,0.771) (0.064,0.185) (0.457,0.669) (0.794,0.053)
    (0.086,0.255) (0.621,0.386) (0.447,0.805) (0.824,0.547) (0.792,0.406) (0.968,0.044)
    (0.714,0.193) (0.549,0.289) (0.105,0.003) (0.205,0.258) (0.466,0.818) (0.118,0.968)
    (0.944,0.245) (0.641,0.360) (0.697,0.087) (0.470,0.588) (0.626,0.693) (0.891,0.241)
    (0.153,0.391) (0.259,0.423) (0.295,0.651) (0.952,0.154) (0.518,0.678) (0.112,0.196)
    (0.066,0.829) (0.963,0.014) (0.908,0.270) (0.715,0.237) (0.076,0.745) (0.042,0.700)
    (0.283,0.487) (0.892,0.243) (0.468,0.355) (0.190,0.474) (0.182,0.799) (0.615,0.283)
    (0.671,0.733) (0.265,0.012) (0.170,0.154) (0.433,0.598) (0.599,0.436) (0.231,0.065)
    (0.169,0.358) (0.088,0.151) (0.045,0.623) (0.493,0.244) (0.526,0.424) (0.353,0.645)
    (0.219,0.940) (0.618,0.159) (0.257,0.159) (0.179,0.307) (0.823,0.457) (0.277,0.419)
    (0.267,0.734) (0.059,0.289) (0.323,0.182) (0.738,0.180) (0.300,0.056) (0.756,0.364)
    (0.656,0.217) (0.022,0.322) (0.349,0.847) (0.050,0.522) (0.257,0.439) (0.222,0.894)
    (0.302,0.945) (0.109,0.430) (0.443,0.709) (0.345,0.492) (0.070,0.311) (0.725,0.106)
    (0.463,0.237) (0.165,0.121) (0.379,0.392) (0.728,0.604) (0.482,0.659) (0.545,0.310)
    (0.454,0.721) (0.553,0.632) (0.776,0.131) (0.817,0.343) (0.368,0.243) (0.056,0.472)
    (0.659,0.223) (0.803,0.368) (0.890,0.326) (0.092,0.730) (0.108,0.870) (0.412,0.780)
    (0.053,0.544) (0.677,0.130) (0.608,0.143) (0.304,0.365) (0.620,0.130) (0.185,0.781)
    (0.716,0.454) (0.282,0.692) (0.468,0.312) (0.875,0.202) (0.646,0.682) (0.921,0.255)
    (0.164,0.896) (0.204,0.222) (0.845,0.517) (0.736,0.629) (0.011,0.195) (0.013,0.907)
    (0.236,0.171) (0.583,0.203) (0.864,0.245) (0.916,0.399) (0.523,0.325) (0.737,0.620)
    (0.689,0.104) (0.563,0.197) (0.053,0.330) (0.746,0.025) (0.049,0.657) (0.882,0.439)
    (0.624,0.257) (0.273,0.475) (0.219,0.454) (0.148,0.924) (0.648,0.094) (0.193,0.829)
    (0.372,0.260) (0.023,0.675) (0.170,0.752) (0.746,0.442) (0.047,0.644) (0.026,0.086)
    (0.409,0.121) (0.452,0.851) (0.174,0.798) (0.998,0.037) (0.423,0.708) (0.732,0.280)
    (0.809,0.371) (0.201,0.518) (0.664,0.083) (0.575,0.344) (0.705,0.045) (0.040,0.015)
    (0.177,0.922) (0.407,0.519) (0.471,0.335) (0.475,0.571) (0.006,0.768) (0.140,0.944)
    (0.044,0.708) (0.769,0.542) (0.096,0.812) (0.321,0.288) (0.096,0.436) (0.673,0.277)
    (0.076,0.155) (0.143,0.456) (0.910,0.112) (0.151,0.212) (0.573,0.430) (0.418,0.145)
    (0.824,0.479) (0.866,0.030) (0.458,0.368) (0.521,0.161) (0.591,0.314) (0.221,0.799)
    (0.782,0.486) (0.242,0.588) (0.656,0.203) (0.145,0.220) (0.002,0.939) (0.174,0.834)
    (0.512,0.057) (0.211,0.258) (0.764,0.368) (0.142,0.946) (0.515,0.067) (0.441,0.146)
    (0.693,0.097) (0.864,0.469) (0.812,0.110) (0.211,0.111) (0.110,0.589) (0.181,0.205)
    (0.758,0.282) (0.346,0.496) (0.249,0.595) (0.187,0.449) (0.602,0.112) (0.770,0.102)
    (0.045,0.786) (0.304,0.495) (0.787,0.609) (0.849,0.395) (0.903,0.398) (0.908,0.309)
    (0.644,0.592) (0.361,0.514) (0.491,0.678) (0.540,0.318) (0.619,0.173) (0.338,0.650)
    (0.239,0.895) (0.541,0.681) (0.664,0.162) (0.146,0.570) (0.528,0.310) (0.375,0.259)
    (0.631,0.429) (0.100,0.227) (0.152,0.629) (0.554,0.393) (0.547,0.629) (0.217,0.056)
    (0.962,0.186) (0.242,0.704) (0.986,0.124) (0.402,0.856) (0.085,0.250) (0.103,0.416)
    (0.592,0.547) (0.256,0.400) (0.101,0.146) (0.294,0.198) (0.617,0.298) (0.014,0.094)
    (0.878,0.109) (0.641,0.419) (0.769,0.509) (0.231,0.211) (0.235,0.234) (0.251,0.299)
    (0.598,0.528) (0.069,0.421) (0.320,0.156) (0.293,0.133) (0.259,0.925) (0.877,0.452)
    (0.550,0.701) (0.701,0.694) (0.736,0.591) (0.133,0.687) (0.669,0.162) (0.195,0.637)
    (0.386,0.811) (0.277,0.666) (0.615,0.111) (0.278,0.353) (0.302,0.091) (0.236,0.537)
    (0.226,0.636) (0.415,0.456) (0.207,0.787) (0.077,0.944) (0.549,0.352) (0.521,0.394)
    (0.156,0.152) (0.416,0.452) (0.804,0.122) (0.454,0.133) (0.254,0.493) (0.877,0.132)
    (0.174,0.706) (0.527,0.088) (0.743,0.329) (0.800,0.568) (0.091,0.362) (0.370,0.709)};
  \addplot[only marks,mark=*,mark size=0.7pt,PlotRed] coordinates {(0.967,0.683) (0.891,0.776) (0.680,0.849) (0.892,0.614) (0.898,0.835) (0.385,0.974)
    (0.504,0.937) (0.927,0.785) (0.982,0.446) (0.910,0.962) (0.903,0.867) (0.977,0.615)
    (0.974,0.605) (0.883,0.559) (0.907,0.671) (0.569,0.961) (0.711,0.738) (0.973,0.268)
    (0.493,0.914) (0.737,0.882) (0.789,0.728) (0.973,0.345) (0.893,0.956) (0.989,0.963)
    (0.988,0.246) (0.659,0.933) (0.779,0.837) (0.814,0.687) (0.916,0.919) (0.487,0.953)
    (0.814,0.918) (0.850,0.769) (0.662,0.847) (0.804,0.618) (0.748,0.881) (0.805,0.966)
    (0.376,0.989) (0.547,0.950) (1.000,0.307) (0.537,0.997) (0.570,0.840) (0.827,0.824)
    (0.612,0.838) (0.657,0.792) (0.892,0.895) (0.884,0.611) (0.932,0.622) (0.925,0.515)
    (0.692,0.972) (0.925,0.603) (0.841,0.692) (0.683,0.896) (0.777,0.967) (0.812,0.635)
    (0.852,0.695) (0.936,0.806) (0.973,0.675) (0.939,0.752) (0.713,0.770) (0.700,0.751)
    (0.917,0.892) (0.929,0.663) (0.776,0.669) (0.855,0.658) (0.822,0.620) (0.948,0.934)
    (0.932,0.534) (0.673,0.947) (0.845,0.861) (0.594,0.980) (0.191,0.987) (0.932,0.846)
    (0.614,0.829) (0.536,0.993) (0.984,0.210) (0.565,0.987) (0.930,0.940) (0.589,0.944)
    (0.370,0.929) (0.681,0.958) (0.738,0.996) (0.946,0.655)};
  \addplot[black,thick,domain=0:90,samples=60] ({cos(x)},{sin(x)});
\end{axis}
\end{tikzpicture}
\caption{Monte Carlo estimation of $\pi$: $318$ of $400$ random points (blue) lie
  inside the quarter circle.}
\label{fig:monte-carlo-pi}
\end{figure}

Monte Carlo also computes integrals: $\int_0^1e^{-x^2}\,\dd x$, which has no
formula in terms of elementary functions, equals $\E\bigl[e^{-U^2}\bigr]$ for $U$
uniform on $[0, 1]$. Averaging $e^{-u^2}$ over $10\,000$ random numbers $u$ gave
$0.7460$, with standard error $0.002$; the true value is $0.7468$. The method
gets more accurate only slowly (the square-root law again), but it works the
same way in a thousand dimensions, where other methods are hopeless.

\begin{mlnote}[Monte Carlo is everywhere in machine learning]
Whenever an expectation is too expensive to compute exactly, machine learning
replaces it by an average over random samples.
\begin{itemize}
  \item A \textbf{mini-batch gradient} is a Monte Carlo estimate of the full
    gradient (\cref{sec:sgd}): unbiased, with a standard error proportional to
    $1/\sqrt B$ for batch size $B$. Quadrupling the batch only halves the noise.
  \item A \textbf{test-set evaluation} is a Monte Carlo estimate of a model's
    expected loss on new data.
  \item Generative models, reinforcement learning and Bayesian methods all rely
    on averages over random samples to approximate expectations that cannot be
    computed exactly.
\end{itemize}
\end{mlnote}

\sectionexercises

\begin{exercise}
Use Chebyshev's inequality to find a number of die rolls $n$ that guarantees
$P\bigl(\abs{\bar X_n - 3.5} \ge 0.1\bigr) \le 0.05$. (The variance of one roll
is $\tfrac{35}{12}$.)
\end{exercise}
\begin{solution}
We need $\tfrac{35/12}{n\times0.1^2} \le 0.05$, i.e.\ $n \ge \tfrac{2.9167}{0.0005}
\approx 5833.3$, so $n = 5834$ rolls suffice. (Chebyshev is pessimistic; the
central limit theorem shows that about $1100$ are enough.)
\end{solution}

\begin{exercise}
The five numbers $0.1, 0.4, 0.5, 0.8, 0.9$ are random samples from the uniform
distribution on $[0, 1]$. Use them to give a Monte Carlo estimate of $\E[X^2]$,
and compare with the exact value.
\end{exercise}
\begin{solution}
$\tfrac15(0.01 + 0.16 + 0.25 + 0.64 + 0.81) = \tfrac{1.87}{5} = 0.374$. The exact
value is $\int_0^1x^2\,\dd x = \tfrac13 \approx 0.333$.
\end{solution}

\begin{exercise}
With a batch size of $32$, the standard deviation of a mini-batch gradient
component is $0.4$. Roughly which batch size would bring it down to $0.1$?
\end{exercise}
\begin{solution}
The standard deviation scales like $1/\sqrt B$. Dividing it by $4$ requires
multiplying $B$ by $16$: $B = 512$.
\end{solution}

\begin{exercise}
A fair coin has landed heads ten times in a row. What is the probability of heads
on the next toss? What was the probability, before any toss, of eleven heads in a
row?
\end{exercise}
\begin{solution}
The next toss is independent of the past: $\tfrac12$. Eleven heads in a row had
probability $\left(\tfrac12\right)^{11} = \tfrac{1}{2048}$ before the first toss;
the rarity of the whole sequence does not change the next toss.
\end{solution}

%----------------------------------------------------------------------
\section{The central limit theorem and confidence intervals}\label{sec:clt}
%----------------------------------------------------------------------

The law of large numbers says that $\bar X_n$ approaches $\mu$. The central limit
theorem says something much more detailed: \emph{how} $\bar X_n$ fluctuates
around $\mu$. The answer is always the same: like a normal distribution.

\begin{theorem}[Central limit theorem]\label{thm:clt}
Let $X_1, X_2, \dots$ be i.i.d.\ with mean $\mu$ and finite variance
$\sigma^2 > 0$. For large $n$, the sample mean is approximately normal:
\[
  \keyeq{\bar X_n \approx \mathcal{N}\!\left(\mu, \frac{\sigma^2}{n}\right)}
\]
Precisely: for every number $z$,
\[
  P\left(\frac{\bar X_n - \mu}{\sigma/\sqrt n} \le z\right) \to P(Z \le z)
  \qquad\text{as } n \to \infty,
\]
where $Z \sim \mathcal{N}(0, 1)$.
\end{theorem}
\index{central limit theorem}

\begin{figure}[htbp]
\centering
\begin{tikzpicture}
\begin{groupplot}[group style={group size=2 by 2,horizontal sep=1.2cm,vertical sep=1.3cm},
  mlplot,axis lines=left,width=6.4cm,height=4.4cm,xmin=0.5,xmax=6.5,ymin=0,
  xtick={1,...,6},yticklabel style={/pgf/number format/fixed},
  title style={font=\small,yshift=-3pt}]
  \nextgroupplot[title={$n = 1$},ymax=0.3,ytick={0,0.1,0.2,0.3}]
    \addplot[ybar interval,fill=MLdef!35,draw=MLdef!70,thin] coordinates {(0.5000,0.1667) (1.5000,0.1667) (2.5000,0.1667) (3.5000,0.1667) (4.5000,0.1667) (5.5000,0.1667) (6.5000,0.1667)};
    \addplot[PlotRed,domain=0.5:6.5,samples=100]{exp(-(x-3.5)^2/(2*35/12))/sqrt(2*pi*35/12)};
  \nextgroupplot[title={$n = 2$},ymax=0.42,ytick={0,0.1,0.2,0.3,0.4}]
    \addplot[ybar interval,fill=MLdef!35,draw=MLdef!70,thin] coordinates {(0.7500,0.0556) (1.2500,0.1111) (1.7500,0.1667) (2.2500,0.2222) (2.7500,0.2778)
      (3.2500,0.3333) (3.7500,0.2778) (4.2500,0.2222) (4.7500,0.1667) (5.2500,0.1111)
      (5.7500,0.0556) (6.2500,0.0556)};
    \addplot[PlotRed,domain=0.5:6.5,samples=100]{exp(-(x-3.5)^2/(2*35/24))/sqrt(2*pi*35/24)};
  \nextgroupplot[title={$n = 5$},ymax=0.65,ytick={0,0.2,0.4,0.6}]
    \addplot[ybar interval,fill=MLdef!35,draw=MLdef!70,thin] coordinates {(0.9000,0.0006) (1.1000,0.0032) (1.3000,0.0096) (1.5000,0.0225) (1.7000,0.0450)
      (1.9000,0.0810) (2.1000,0.1318) (2.3000,0.1961) (2.5000,0.2701) (2.7000,0.3472)
      (2.9000,0.4186) (3.1000,0.4726) (3.3000,0.5015) (3.5000,0.5015) (3.7000,0.4726)
      (3.9000,0.4186) (4.1000,0.3472) (4.3000,0.2701) (4.5000,0.1961) (4.7000,0.1318)
      (4.9000,0.0810) (5.1000,0.0450) (5.3000,0.0225) (5.5000,0.0096) (5.7000,0.0032)
      (5.9000,0.0006) (6.1000,0.0006)};
    \addplot[PlotRed,domain=0.5:6.5,samples=100]{exp(-(x-3.5)^2/(2*35/60))/sqrt(2*pi*35/60)};
  \nextgroupplot[title={$n = 25$},ymax=1.35,ytick={0,0.5,1}]
    \addplot[ybar interval,fill=MLdef!35,draw=none] coordinates {(2.1800,0.0006) (2.2200,0.0010) (2.2600,0.0016) (2.3000,0.0025) (2.3400,0.0039)
      (2.3800,0.0058) (2.4200,0.0087) (2.4600,0.0127) (2.5000,0.0181) (2.5400,0.0256)
      (2.5800,0.0355) (2.6200,0.0485) (2.6600,0.0652) (2.7000,0.0863) (2.7400,0.1125)
      (2.7800,0.1445) (2.8200,0.1827) (2.8600,0.2278) (2.9000,0.2798) (2.9400,0.3389)
      (2.9800,0.4046) (3.0200,0.4763) (3.0600,0.5528) (3.1000,0.6329) (3.1400,0.7146)
      (3.1800,0.7959) (3.2200,0.8744) (3.2600,0.9478) (3.3000,1.0135) (3.3400,1.0693)
      (3.3800,1.1131) (3.4200,1.1432) (3.4600,1.1586) (3.5000,1.1586) (3.5400,1.1432)
      (3.5800,1.1131) (3.6200,1.0693) (3.6600,1.0135) (3.7000,0.9478) (3.7400,0.8744)
      (3.7800,0.7959) (3.8200,0.7146) (3.8600,0.6329) (3.9000,0.5528) (3.9400,0.4763)
      (3.9800,0.4046) (4.0200,0.3389) (4.0600,0.2798) (4.1000,0.2278) (4.1400,0.1827)
      (4.1800,0.1445) (4.2200,0.1125) (4.2600,0.0863) (4.3000,0.0652) (4.3400,0.0485)
      (4.3800,0.0355) (4.4200,0.0256) (4.4600,0.0181) (4.5000,0.0127) (4.5400,0.0087)
      (4.5800,0.0058) (4.6200,0.0039) (4.6600,0.0025) (4.7000,0.0016) (4.7400,0.0010)
      (4.7800,0.0006) (4.8200,0.0006)};
    \addplot[PlotRed,domain=0.5:6.5,samples=150]{exp(-(x-3.5)^2/(2*35/300))/sqrt(2*pi*35/300)};
\end{groupplot}
\end{tikzpicture}
\caption{The central limit theorem in action. The bars show the exact
  distribution of the average of $n$ rolls of a fair die, drawn as a density
  (bar heights are probabilities divided by the bar width $\frac1n$). The red
  curves are the normal densities $\mathcal{N}(3.5, \frac{35}{12n})$ predicted by
  the theorem. A flat distribution turns into a bell curve that becomes narrower
  as $n$ grows.}
\label{fig:clt}
\end{figure}

\begin{intuition}[Why averages become bell-shaped]
An average of many independent pieces has many ways to come out near the middle
and few ways to come out at the extremes: to average $6$ over $25$ dice, every
single die must show a $6$. Remarkably, the limiting shape does not depend on the
shape of the individual distribution (flat, skewed, two-humped): only its mean
and variance survive. This is why the normal distribution appears whenever many
small independent effects add up.
\end{intuition}

\subsection{Confidence intervals}

The central limit theorem lets us attach error bars to an estimate. Since
$\bar X$ is approximately normal with mean $\mu$ and standard deviation
$\operatorname{SE} = \sigma/\sqrt n$, and a normal random variable lies within
$1.96$ standard deviations of its mean with probability $0.95$,
\[
  P\bigl(\bar X - 1.96\operatorname{SE} \le \mu \le \bar X + 1.96\operatorname{SE}\bigr)
  \approx 0.95 .
\]

\begin{definition}[Confidence interval]\label{def:confidence-interval}
An approximate $95\%$ \term{confidence interval} for a mean $\mu$ is
\[
  \bar x \pm 1.96\,\frac{s}{\sqrt n}
  = \left[\bar x - 1.96\,\frac{s}{\sqrt n},\ \bar x + 1.96\,\frac{s}{\sqrt n}\right].
\]
\end{definition}

\begin{example}[Error bars for an accuracy]\label{ex:accuracy-ci}
For the classifier of \cref{ex:accuracy-se}, with accuracy $0.85$ on $1000$ test
examples, the standard error is $0.0113$, and an approximate $95\%$ confidence
interval for the true accuracy is
\[
  0.85 \pm 1.96\times0.0113 = 0.85 \pm 0.022 = [0.828,\ 0.872].
\]
\end{example}

\begin{pitfall}[What ``$95\%$ confident'' means]
The true accuracy is a fixed number; it is the interval that is random. The
$95\%$ describes the \emph{procedure}: if we repeated the whole experiment many
times, with new test sets, about $95\%$ of the intervals computed this way would
contain the true value. It does not mean that the true value lies in our
particular interval with probability $95\%$ (in the Bayesian approach of
\cref{sec:map} one can make such statements, using a prior).
\end{pitfall}

\begin{mlnote}[Error bars in machine learning]
Results computed from finite samples should come with error bars: report the
mean and standard error over several training runs with different random seeds,
or a confidence interval for a test metric. The central limit theorem also
explains why the input $\sum_iw_ix_i$ to a neuron with many inputs is
approximately normal, which is the starting point of the analysis of weight
initialisation (\cref{sec:variance}).
\end{mlnote}

\sectionexercises

\begin{exercise}
A die is rolled $100$ times. What is the approximate distribution of the average
roll? Estimate the probability that the average exceeds $3.84$.
\end{exercise}
\begin{solution}
Approximately $\mathcal{N}\bigl(3.5, \tfrac{35/12}{100}\bigr)$, with standard
deviation $\sqrt{0.02917} \approx 0.171$. Since $3.84$ is $\tfrac{0.34}{0.171}
\approx 2$ standard deviations above the mean, the probability is about
$2.5\%$ (more precisely $0.023$).
\end{solution}

\begin{exercise}
A model is correct on $360$ of $400$ test examples. Give an approximate $95\%$
confidence interval for its accuracy.
\end{exercise}
\begin{solution}
$\hat p = 0.9$ and $\operatorname{SE} = \sqrt{\tfrac{0.9\times0.1}{400}} = 0.015$, so
the interval is $0.9 \pm 0.029 = [0.871, 0.929]$.
\end{solution}

\begin{exercise}
About how many test examples are needed to measure an accuracy of about $90\%$
to within $\pm1$ percentage point (as a $95\%$ confidence interval)?
\end{exercise}
\begin{solution}
We need $1.96\sqrt{\tfrac{0.9\times0.1}{n}} \le 0.01$, i.e.\ $n \ge
\tfrac{1.96^2\times0.09}{0.0001} \approx 3457$: about $3500$ examples.
\end{solution}

\begin{exercise}
Let $U_1, \dots, U_{12}$ be independent and uniform on $[0, 1]$ (mean
$\tfrac12$, variance $\tfrac1{12}$). Find the mean and variance of
$S = U_1 + \dots + U_{12} - 6$. What does the central limit theorem suggest about
its distribution?
\end{exercise}
\begin{solution}
$\E[S] = 12\times\tfrac12 - 6 = 0$ and $\Var(S) = 12\times\tfrac1{12} = 1$. As a sum
of many independent pieces, $S$ is approximately standard normal; this was once a
popular way of generating normal random numbers.
\end{solution}

%----------------------------------------------------------------------
\section{Maximum likelihood estimation}\label{sec:mle}
%----------------------------------------------------------------------

Most models in machine learning are \emph{parametric}: a family of
distributions $p(x \mid \theta)$, one for each value of a parameter $\theta$
(which may be a vector of millions of weights). Training means choosing
$\theta$. The most important principle for doing this is maximum likelihood.

\begin{picturethis}[Which coin produced these tosses?]
A coin is tossed $10$ times and shows $7$ heads and $3$ tails. Suppose the
probability of heads is $\theta$. How probable was this particular sequence of
tosses?
\begin{center}
\begin{tabular}{lccc}
\toprule
$\theta$ & $0.5$ & $0.7$ & $0.9$\\
\midrule
$P(\text{data} \mid \theta) = \theta^7(1 - \theta)^3$ & $0.00098$ & $0.00222$ &
  $0.00048$\\
\bottomrule
\end{tabular}
\end{center}
Of these values, $\theta = 0.7$ makes the observed data most probable, more than
twice as probable as a fair coin does. Maximum likelihood says: \emph{choose the
parameter value under which the observed data would have been most probable}.
\end{picturethis}

\begin{definition}[Likelihood]\label{def:likelihood}
For data $x_1, \dots, x_n$ and a model $p(x \mid \theta)$, the
\term[likelihood]{likelihood} is the probability (or density) of the data,
viewed as a function of the parameter:
\[
  \mathcal{L}(\theta) = p(x_1, \dots, x_n \mid \theta)
  = \prod_{i=1}^n p(x_i \mid \theta) ,
\]
where the product form holds for i.i.d.\ data. Its logarithm
$\ln\mathcal{L}(\theta) = \sum_{i=1}^n\ln p(x_i \mid \theta)$ is the
\term{log-likelihood}.
\end{definition}

\begin{definition}[Maximum likelihood estimate]\label{def:mle}
The \term[maximum likelihood estimation]{maximum likelihood estimate} (MLE) is
the parameter value that maximises the likelihood:
\[
  \keyeq{\hat\theta_{\mathrm{MLE}} = \argmax_\theta\,\ln\mathcal{L}(\theta)
  = \argmin_\theta\,\Bigl(-\sum_{i=1}^n\ln p(x_i \mid \theta)\Bigr)}
\]
The quantity being minimised on the right is the
\term{negative log-likelihood} (NLL).
\end{definition}

Maximising $\ln\mathcal{L}$ instead of $\mathcal{L}$ gives the same answer,
because $\ln$ is increasing, and it is far more convenient: sums are easier to
differentiate than products, and a product of thousands of probabilities would
underflow to $0$ on a computer (\cref{sec:logarithms}).

\begin{example}[Maximum likelihood for a coin]\label{ex:mle-coin}
Find the MLE of $\theta = P(\text{heads})$ after $h$ heads in $n$ tosses; in
particular for $7$ heads in $10$ tosses.

\sol The log-likelihood is $\ln\mathcal{L}(\theta) = h\ln\theta + (n - h)\ln(1 -
\theta)$. Setting its derivative to zero,
\[
  \frac{h}{\theta} - \frac{n - h}{1 - \theta} = 0
  \quad\Longrightarrow\quad h(1 - \theta) = (n - h)\,\theta
  \quad\Longrightarrow\quad \hat\theta = \frac hn .
\]
The second derivative $-\frac{h}{\theta^2} - \frac{n-h}{(1-\theta)^2}$ is negative,
so this is a maximum. For $7$ heads in $10$ tosses, $\hat\theta = 0.7$
(\cref{fig:likelihood}): the MLE is the observed frequency, as common sense
suggests.
\end{example}

\begin{figure}[htbp]
\centering
\begin{tikzpicture}
\begin{groupplot}[group style={group size=2 by 1,horizontal sep=1.6cm},
  mlplot,axis lines=left,width=6.6cm,height=5cm,xmin=0,xmax=1,
  xlabel={$\theta$},
  xlabel style={at={(ticklabel cs:0.5)},anchor=north},
  ylabel style={at={(ticklabel cs:0.5)},anchor=south,rotate=90}]
  \nextgroupplot[ymin=0,ymax=2.6,ylabel={$\mathcal{L}(\theta)\times10^{3}$}]
    \addplot[PlotBlue,domain=0:1,samples=150]{1000*x^7*(1-x)^3};
    \draw[dashed,gray] (axis cs:0.7,0) -- (axis cs:0.7,2.2236);
    \fill[PlotRed] (axis cs:0.7,2.2236) circle (2pt);
  \nextgroupplot[ymin=-22,ymax=-4,ylabel={$\ln\mathcal{L}(\theta)$}]
    \addplot[PlotBlue,domain=0.05:0.99,samples=150]{7*ln(x)+3*ln(1-x)};
    \draw[dashed,gray] (axis cs:0.7,-22) -- (axis cs:0.7,-6.109);
    \fill[PlotRed] (axis cs:0.7,-6.109) circle (2pt);
\end{groupplot}
\end{tikzpicture}
\caption{The likelihood (left) and log-likelihood (right) of $\theta$ after
  $7$ heads in $10$ tosses. Both are maximised at $\hat\theta = 0.7$; the
  logarithm changes the shape of the curve but not the location of its maximum.}
\label{fig:likelihood}
\end{figure}

\begin{example}[Maximum likelihood for a normal distribution]\label{ex:mle-gaussian}
Data $x_1, \dots, x_n$ are modelled as i.i.d.\ from $\mathcal{N}(\mu, \sigma^2)$.
Find the MLEs of $\mu$ and $\sigma^2$.

\sol Taking the logarithm of the normal density (\cref{sec:distributions}) and
adding over the data,
\[
  \ln\mathcal{L}(\mu, \sigma^2) = -\frac n2\ln(2\pi\sigma^2)
  - \frac{1}{2\sigma^2}\sum_{i=1}^n(x_i - \mu)^2 .
\]
For any fixed $\sigma^2$, this is largest when $\sum_i(x_i - \mu)^2$ is smallest,
which happens at $\hat\mu = \bar x$ (\cref{ex:mean-minimises}). Next, writing
$v = \sigma^2$ and differentiating with respect to $v$,
\[
  -\frac{n}{2v} + \frac{1}{2v^2}\sum_{i=1}^n(x_i - \bar x)^2 = 0
  \quad\Longrightarrow\quad
  \hat\sigma^2 = \frac1n\sum_{i=1}^n(x_i - \bar x)^2 .
\]
Note that the MLE of the variance divides by $n$, so it is slightly biased
(\cref{thm:bessel}). Maximum likelihood estimates need not be unbiased.
\end{example}

For a categorical distribution with $K$ outcomes, observed $n_1, \dots, n_K$
times in $N$ trials, the MLE is $\hat p_k = \frac{n_k}{N}$, the observed
frequencies; we derived this with a Lagrange multiplier in
\cref{ex:lagrange-categorical}.

\begin{remark}[How good is maximum likelihood?]
If the model is correct, then as the amount of data grows the MLE converges to
the true parameter, its sampling distribution becomes approximately normal, and
no other (approximately unbiased) estimator has a smaller variance. With little
data, however, it can be badly overconfident: after $3$ heads in $3$ tosses the
MLE is $\hat\theta = 1$, which declares tails impossible. Priors
(\cref{sec:map}) fix this.
\end{remark}

\begin{mlnote}[Training is maximum likelihood]
Most of machine learning is maximum likelihood in disguise. A language model is
trained by maximising the log-likelihood of its training text: by the chain rule
of probability (\cref{sec:conditional}), this is the sum of the log-probabilities
that it assigns to each actual next token. Minimising the average negative
log-likelihood with stochastic gradient descent is the training procedure of
nearly every modern neural network, as the next section shows.
\end{mlnote}

\sectionexercises

\begin{exercise}
A coin shows $12$ heads in $40$ tosses. What is the maximum likelihood estimate
of the probability of heads?
\end{exercise}
\begin{solution}
$\hat\theta = \tfrac{12}{40} = 0.3$.
\end{solution}

\begin{exercise}
Counts $x_1, \dots, x_n$ are modelled as Poisson with
$p(x \mid \lambda) = \frac{\lambda^xe^{-\lambda}}{x!}$. Show that the MLE is
$\hat\lambda = \bar x$, and compute it for the data $2, 0, 3, 1, 4$.
\end{exercise}
\begin{solution}
$\ln\mathcal{L}(\lambda) = \sum_i(x_i\ln\lambda - \lambda - \ln x_i!)$. Its
derivative $\tfrac{1}{\lambda}\sum_ix_i - n$ vanishes at $\hat\lambda =
\tfrac1n\sum_ix_i = \bar x$ (and the second derivative $-\sum_ix_i/\lambda^2$ is
negative). For the data, $\hat\lambda = \tfrac{10}{5} = 2$.
\end{solution}

\begin{exercise}
Waiting times $0.5, 1.5, 2.0, 1.0$ are modelled as exponential with density
$\lambda e^{-\lambda t}$. Find the MLE of $\lambda$.
\end{exercise}
\begin{solution}
$\ln\mathcal{L}(\lambda) = n\ln\lambda - \lambda\sum_it_i$, with derivative
$\tfrac n\lambda - \sum_it_i = 0$, so $\hat\lambda = \tfrac{n}{\sum_it_i} =
\tfrac{1}{\bar t} = \tfrac{1}{1.25} = 0.8$.
\end{solution}

\begin{exercise}
Find the MLEs of $\mu$ and $\sigma^2$ for the data of \cref{ex:summary-stats},
assuming a normal model, and compare $\hat\sigma^2$ with $s^2$.
\end{exercise}
\begin{solution}
$\hat\mu = \bar x = 5$ and $\hat\sigma^2 = \tfrac{32}{8} = 4$, slightly smaller
than $s^2 = \tfrac{32}{7} \approx 4.57$.
\end{solution}

\begin{exercise}
The log-likelihood of a data set is $-120.5$ under model A and $-118.2$ under
model B. Which model explains the data better, and how many times larger is its
likelihood?
\end{exercise}
\begin{solution}
Model B, whose log-likelihood is larger. The ratio of the likelihoods is
$e^{-118.2}/e^{-120.5} = e^{2.3} \approx 10$.
\end{solution}

%----------------------------------------------------------------------
\section{Loss functions from likelihoods}\label{sec:losses-from-likelihood}
%----------------------------------------------------------------------

In supervised learning the model describes the \emph{conditional} distribution
$p(y \mid \vx, \vtheta)$ of a label $y$ given an input $\vx$. Maximum likelihood
then says: minimise $-\sum_i\ln p(y_i \mid \vx_i, \vtheta)$. Choosing a
distribution for the labels determines the loss function.

\subsection{Regression: Gaussian noise gives squared error}

Suppose that the targets are a deterministic function plus Gaussian noise,
\[
  y = f(\vx; \vtheta) + \varepsilon,
  \qquad \varepsilon \sim \mathcal{N}(0, \sigma^2),
\]
so that $p(y \mid \vx, \vtheta)$ is the normal density with mean $f(\vx;
\vtheta)$ and variance $\sigma^2$ (\cref{fig:gaussian-noise}).

\begin{figure}[htbp]
\centering
\begin{tikzpicture}
\begin{axis}[mlplot,axis lines=left,width=10cm,height=6.5cm,xmin=0,xmax=6.3,ymin=0,ymax=5.2,
  xlabel={$x$},ylabel={$y$},
  xlabel style={at={(ticklabel cs:0.5)},anchor=north},
  ylabel style={at={(ticklabel cs:0.5)},anchor=south,rotate=90}]
  \addplot[only marks,mark=*,mark size=1.4pt,PlotBlue] coordinates {(0.40,0.75) (0.54,1.28) (0.56,1.27) (0.59,1.19) (0.63,0.95) (0.65,1.18) (1.53,1.38)
    (1.57,1.31) (1.77,1.96) (1.84,1.53) (2.37,2.59) (2.42,2.46) (2.51,1.55) (2.65,2.42)
    (2.96,2.10) (3.08,2.55) (3.30,2.67) (3.82,2.22) (3.95,2.90) (3.97,2.89) (4.65,3.66)
    (4.66,2.92) (4.86,3.69) (5.00,3.11) (5.01,3.39) (5.05,3.23) (5.14,4.07) (5.15,3.77)
    (5.56,4.63) (5.70,4.07)};
  \addplot[MLdef,very thick,domain=0:6.3]{0.5*x+1};
  \addplot[PlotRed,thick,domain=-3:3,samples=60,variable=\t] ({1.5+0.9*exp(-(\t)^2/2)},{1.75+0.35*\t});
  \addplot[PlotRed,thick,domain=-3:3,samples=60,variable=\t] ({3+0.9*exp(-(\t)^2/2)},{2.5+0.35*\t});
  \addplot[PlotRed,thick,domain=-3:3,samples=60,variable=\t] ({4.5+0.9*exp(-(\t)^2/2)},{3.25+0.35*\t});
  \draw[PlotRed] (axis cs:1.5,0.7) -- (axis cs:1.5,2.8);
  \draw[PlotRed] (axis cs:3,1.45) -- (axis cs:3,3.55);
  \draw[PlotRed] (axis cs:4.5,2.2) -- (axis cs:4.5,4.3);
  \node[MLdef,font=\small,anchor=north east] at (axis cs:6.25,3.75) {$f(x)$};
\end{axis}
\end{tikzpicture}
\caption{The probabilistic view of regression. For each input $x$, the target
  $y$ is drawn from a normal distribution (red) centred on the line $f(x)$. Maximum
  likelihood for this model is least squares.}
\label{fig:gaussian-noise}
\end{figure}

The negative log-likelihood is
\[
  -\ln\mathcal{L}(\vtheta)
  = \frac{1}{2\sigma^2}\sum_{i=1}^n\bigl(y_i - f(\vx_i; \vtheta)\bigr)^2
  + n\ln\bigl(\sigma\sqrt{2\pi}\bigr) .
\]
The second term does not depend on $\vtheta$, and the first is a positive
multiple of the sum of squared errors. Therefore:

\begin{theorem}[Least squares is maximum likelihood]\label{thm:ls-mle}
If the noise is Gaussian with constant variance, maximum likelihood estimation
of $\vtheta$ is the same as minimising the sum of squared errors (equivalently,
the mean squared error) $\sum_i\bigl(y_i - f(\vx_i; \vtheta)\bigr)^2$.
\end{theorem}

This is the statistical reason why squared error is the default loss for
regression. It also explains its weakness: the normal distribution has very thin
tails, so the model regards a large residual as nearly impossible and bends
itself to reduce it, which makes least squares sensitive to outliers. If the
noise is instead modelled by the heavier-tailed \emph{Laplace distribution},
with density $\frac{1}{2b}e^{-\abs\varepsilon/b}$, maximum likelihood gives the
more robust absolute error (an exercise below).

\subsection{Classification: Bernoulli labels give cross-entropy}

For binary classification, $y \in \{0, 1\}$, and the model outputs the
probability $q = P(y = 1 \mid \vx, \vtheta)$, for example
$q = \sigma(\vw^\top\vx + b)$ with the sigmoid function of \cref{ch:functions};
this model is called \term{logistic regression}. The Bernoulli probability of
the observed label can be written in one formula as $q^y(1 - q)^{1-y}$ (which is
$q$ if $y = 1$ and $1 - q$ if $y = 0$). So the negative log-likelihood is
\[
  \keyeq{-\ln\mathcal{L}(\vtheta)
  = -\sum_{i=1}^n\Bigl[y_i\ln q_i + (1 - y_i)\ln(1 - q_i)\Bigr]}
\]
where $q_i$ is the predicted probability for example $i$. This loss is the
\term{binary cross-entropy} or \emph{log loss}.

\begin{example}[Computing the binary cross-entropy]\label{ex:bce}
Three examples have labels $1, 0, 1$, and a model predicts
$q = 0.9, 0.2, 0.6$. Compute the loss.

\sol For each example the loss is $-\ln$ of the probability given to the true
label: $-\ln0.9 \approx 0.105$, $-\ln(1 - 0.2) = -\ln0.8 \approx 0.223$ and
$-\ln0.6 \approx 0.511$. The total is $0.839$ and the average $0.280$.
\end{example}

\begin{intuition}[Cross-entropy punishes confident mistakes]
The loss on an example is $-\ln(\text{probability assigned to the truth})$. It
is $0$ for a confident correct prediction, about $0.69$ for an undecided
prediction of $0.5$, and it grows without limit as the probability of the truth
goes to $0$: predicting $q = 0.99$ for an example whose label is $0$ costs
$-\ln0.01 \approx 4.6$. A model trained with cross-entropy learns to be confident
only when it has reason to be.
\end{intuition}

For $K$ classes the labels are categorical, the model outputs probabilities
$q_1, \dots, q_K$ (usually through a softmax function), and the negative
log-likelihood is $-\sum_i\ln q_{i, y_i}$, the \emph{cross-entropy} loss. Its
connection with information theory, and its strikingly simple gradient, are the
subject of \cref{ch:dl}.

\begin{keytakeaways}[Choosing a loss means choosing a noise model]
\renewcommand{\arraystretch}{1.3}
\begin{tabular}{@{}lll@{}}
\textbf{Target} & \textbf{Model for $y$ given $\vx$} & \textbf{Negative log-likelihood}\\
\hline
real number & normal, mean $f(\vx)$ & squared error\\
real number & Laplace, centre $f(\vx)$ & absolute error\\
binary label & Bernoulli$(q(\vx))$ & binary cross-entropy\\
one of $K$ classes & categorical$(\vq(\vx))$ & cross-entropy\\
count & Poisson$(\lambda(\vx))$ & $\lambda(\vx) - y\ln\lambda(\vx)$ (Poisson loss)
\end{tabular}
\end{keytakeaways}

\sectionexercises

\begin{exercise}
Four examples have labels $0, 1, 1, 0$, and a model predicts the probabilities
$q = 0.1, 0.7, 0.4, 0.3$ that the label is $1$. Compute the average binary
cross-entropy. Which example contributes most?
\end{exercise}
\begin{solution}
The losses are $-\ln0.9 \approx 0.105$, $-\ln0.7 \approx 0.357$,
$-\ln0.4 \approx 0.916$ and $-\ln0.7 \approx 0.357$, with total $1.735$ and average
$0.434$. The third example, a positive example given only $0.4$, contributes
most.
\end{solution}

\begin{exercise}
Suppose the noise is Laplace: $p(y \mid \vx) = \frac{1}{2b}\exp\bigl(-\abs{y -
f(\vx)}/b\bigr)$. Show that maximum likelihood minimises the sum of absolute errors
$\sum_i\abs{y_i - f(\vx_i)}$.
\end{exercise}
\begin{solution}
$-\ln\mathcal{L} = \sum_i\left(\ln(2b) + \tfrac1b\abs{y_i - f(\vx_i)}\right)
= n\ln(2b) + \tfrac1b\sum_i\abs{y_i - f(\vx_i)}$. The first term and the factor
$\tfrac1b$ do not depend on the model, so minimising the NLL means minimising
$\sum_i\abs{y_i - f(\vx_i)}$.
\end{solution}

\begin{exercise}
A ``model'' predicts the same probability $q$ for every example of a data set
with $h$ positive and $n - h$ negative labels. Which $q$ minimises the binary
cross-entropy?
\end{exercise}
\begin{solution}
The loss is $-h\ln q - (n - h)\ln(1 - q)$, the negative log-likelihood of the coin
in \cref{ex:mle-coin}, so it is minimised by $q = \tfrac hn$, the fraction of
positive labels.
\end{solution}

\begin{exercise}
For a label $y = 1$, compare the binary cross-entropy of the predictions
$q = 0.9$, $q = 0.5$, $q = 0.1$ and $q = 0.001$.
\end{exercise}
\begin{solution}
$-\ln0.9 \approx 0.105$, $-\ln0.5 \approx 0.693$, $-\ln0.1 \approx 2.303$ and
$-\ln0.001 \approx 6.908$. The loss grows slowly at first and then very fast as the
prediction becomes confidently wrong.
\end{solution}

%----------------------------------------------------------------------
\section{Bayesian inference and MAP estimation}\label{sec:map}
%----------------------------------------------------------------------

Maximum likelihood uses only the data. Often we also have prior knowledge: coins
are usually close to fair; the weights of a good model are usually not huge.
Bayes' theorem (\cref{sec:bayes}) combines prior knowledge with data. Treating
the parameter $\theta$ itself as a random variable with a \term{prior}
distribution $p(\theta)$,
\[
  \underbrace{p(\theta \mid \mathcal{D})}_{\text{posterior}}
  = \frac{p(\mathcal{D} \mid \theta)\,p(\theta)}{p(\mathcal{D})}
  \;\propto\; \underbrace{p(\mathcal{D} \mid \theta)}_{\text{likelihood}}\,
  \underbrace{p(\theta)}_{\text{prior}} ,
\]
where $\mathcal{D}$ denotes the data and $\propto$ means ``is proportional to''
(the denominator does not depend on $\theta$).

\begin{definition}[MAP estimate]\label{def:map}
The \term[MAP estimation]{maximum a posteriori} (MAP) estimate is the parameter
value with the highest posterior density:
\[
  \hat\theta_{\mathrm{MAP}} = \argmax_\theta\, p(\theta \mid \mathcal{D})
  = \argmin_\theta\Bigl(\underbrace{-\ln p(\mathcal{D} \mid \theta)}_{\text{NLL}}
  \;\underbrace{-\ln p(\theta)}_{\text{penalty}}\Bigr).
\]
\end{definition}

So MAP estimation is maximum likelihood plus a penalty term coming from the prior.

\begin{example}[A coin with a prior]\label{ex:map-coin}
Take the prior density $p(\theta) = 6\theta(1 - \theta)$ on $[0, 1]$, which
favours values near $\frac12$ without ruling anything out. (It is a valid
density: $\int_0^16\theta(1 - \theta)\,\dd\theta = 6\left(\frac12 -
\frac13\right) = 1$.) After $h$ heads and $t$ tails, find the MAP estimate.

\sol The posterior is proportional to
$\theta^h(1 - \theta)^t\cdot\theta(1 - \theta) = \theta^{h+1}(1 - \theta)^{t+1}$.
This has the same form as the likelihood in \cref{ex:mle-coin}, with $h + 1$
``heads'' and $t + 1$ ``tails'', so it is maximised at
\[
  \hat\theta_{\mathrm{MAP}} = \frac{h + 1}{n + 2}, \qquad n = h + t .
\]
After $3$ heads in $3$ tosses, the MLE is $1$ but the MAP estimate is
$\frac45 = 0.8$ (\cref{fig:prior-posterior}). After $700$ heads in $1000$ tosses,
the two are almost identical ($0.7$ and $0.6996$): with lots of data, the prior's
influence fades away.
\end{example}

\begin{figure}[htbp]
\centering
\begin{tikzpicture}
\begin{axis}[mlplot,axis lines=left,width=10cm,height=5.8cm,xmin=0,xmax=1,ymin=0,ymax=4.3,
  xlabel={$\theta$},ylabel={density},
  xlabel style={at={(ticklabel cs:0.5)},anchor=north},
  ylabel style={at={(ticklabel cs:0.5)},anchor=south,rotate=90},
  legend style={font=\small,at={(0.03,0.97)},anchor=north west,draw=gray!40},
  legend cell align=left]
  \addplot[PlotGreen,dashed,domain=0:1,samples=100]{6*x*(1-x)};
  \addlegendentry{prior $6\theta(1-\theta)$}
  \addplot[PlotGray,densely dotted,very thick,domain=0:1,samples=100]{4*x^3};
  \addlegendentry{likelihood $\theta^3$ (rescaled)}
  \addplot[PlotBlue,very thick,domain=0:1,samples=150]{30*x^4*(1-x)};
  \addlegendentry{posterior $30\theta^4(1-\theta)$}
  \draw[dashed,gray] (axis cs:0.8,0) -- (axis cs:0.8,2.62);
  \node[font=\small,anchor=south] at (axis cs:0.8,2.62) {MAP};
\end{axis}
\end{tikzpicture}
\caption{Bayesian updating for a coin that shows $3$ heads in $3$ tosses. The
  likelihood alone is maximised at $\theta = 1$; combined with the prior, the
  posterior peaks at $\theta = 0.8$.}
\label{fig:prior-posterior}
\end{figure}

The estimate $\frac{h + 1}{n + 2}$ is known as \term{Laplace smoothing} or
\emph{add-one smoothing}. In a naive Bayes spam filter (\cref{ex:naive-bayes}) it
prevents a word that never appeared in the spam training emails from getting
probability $0$, which would make any email containing it impossible to classify
as spam, whatever its other words.

\subsection{Regularisation is a prior}

Consider linear regression with Gaussian noise of variance $\sigma^2$, and give
each weight a normal prior, $w_j \sim \mathcal{N}(0, \tau^2)$, independently:
we believe that weights are usually small. Then
$-\ln p(\vw) = \frac{1}{2\tau^2}\norm{\vw}^2 + \text{constant}$, and the MAP
objective becomes
\[
  \frac{1}{2\sigma^2}\sum_{i=1}^n(y_i - \vw^\top\vx_i)^2
  + \frac{1}{2\tau^2}\norm{\vw}^2 .
\]
Multiplying by $2\sigma^2$ (which does not change the minimiser) gives
\[
  \sum_{i=1}^n(y_i - \vw^\top\vx_i)^2 + \lambda\norm{\vw}^2,
  \qquad \lambda = \frac{\sigma^2}{\tau^2} ,
\]
which is exactly ridge regression (\cref{sec:analytic-optimisation}).

\begin{theorem}[Priors and penalties]\label{thm:priors-penalties}
MAP estimation with a Gaussian prior on the weights is $L^2$ regularisation
(ridge regression, weight decay). With a Laplace prior,
$p(w_j) \propto e^{-\abs{w_j}/b}$, it is $L^1$ regularisation (the lasso), which
favours sparse weight vectors (\cref{sec:norms}).
\end{theorem}

The stronger the prior belief that the weights are small (the smaller $\tau$),
the larger $\lambda$. A flat prior ($\tau \to \infty$) gives $\lambda \to 0$: MAP
becomes maximum likelihood.

\begin{remark}[Full Bayesian inference]
MAP keeps only the peak of the posterior. A fully Bayesian approach keeps the
whole posterior and averages the predictions of all parameter values, weighted
by their posterior probabilities:
$p(y \mid \vx, \mathcal{D}) = \int p(y \mid \vx, \vtheta)\,p(\vtheta \mid
\mathcal{D})\,\dd\vtheta$. This gives honest uncertainty estimates, but the
integral is usually intractable for neural networks; approximations include
ensembles of independently trained networks, Monte Carlo dropout and
variational inference.
\end{remark}

\sectionexercises

\begin{exercise}
With the prior $6\theta(1 - \theta)$ of \cref{ex:map-coin}, find the MAP estimate
after $7$ heads in $10$ tosses, and compare it with the MLE.
\end{exercise}
\begin{solution}
$\hat\theta_{\mathrm{MAP}} = \tfrac{7 + 1}{10 + 2} = \tfrac23 \approx 0.667$, pulled
slightly from the MLE $0.7$ towards $\tfrac12$.
\end{solution}

\begin{exercise}
In ridge regression derived as MAP estimation, the noise variance is
$\sigma^2 = 1$ and the prior variance of each weight is $\tau^2 = 0.1$. What is
$\lambda$? What happens to $\lambda$ as $\tau^2$ grows?
\end{exercise}
\begin{solution}
$\lambda = \tfrac{\sigma^2}{\tau^2} = 10$. As $\tau^2 \to \infty$ (a vaguer and
vaguer prior), $\lambda \to 0$ and we recover least squares.
\end{solution}

\begin{exercise}
Show that with the uniform prior $p(\theta) = 1$ on $[0, 1]$, the MAP estimate
equals the MLE.
\end{exercise}
\begin{solution}
The posterior is proportional to $\mathcal{L}(\theta)\cdot1$, so it has the same
maximiser as the likelihood. ($-\ln p(\theta) = 0$: no penalty.)
\end{solution}

\begin{exercise}
In $50$ spam training emails a certain word never occurs. Estimate the probability
that a spam email contains the word, by maximum likelihood and with Laplace
smoothing. Why is the first answer dangerous for a naive Bayes classifier?
\end{exercise}
\begin{solution}
MLE: $\tfrac{0}{50} = 0$. Laplace: $\tfrac{0 + 1}{50 + 2} = \tfrac{1}{52} \approx
0.019$. With the MLE, every email containing the word gets a spam likelihood of
exactly $0$, so it can never be classified as spam, no matter how ``spammy'' its
other words are.
\end{solution}

%----------------------------------------------------------------------
\section{Hypothesis testing and p-values}\label{sec:hypothesis-testing}
%----------------------------------------------------------------------

A new model scores better than the old one on a test set. Is it really better,
or was it lucky? Hypothesis testing is a framework for answering such questions.

\begin{definition}[Hypothesis test]\label{def:hypothesis-test}
A hypothesis test has the following ingredients:
\begin{itemize}
  \item a \term{null hypothesis} $H_0$, the sceptical default (``the new model
    is no better'', ``the coin is fair''), and an \term{alternative hypothesis}
    $H_1$;
  \item a \term{test statistic}: a number computed from the data whose
    distribution is known if $H_0$ is true;
  \item the \term{p-value}: the probability, \emph{assuming $H_0$ is true}, of
    obtaining a test statistic at least as extreme as the one observed;
  \item a \term{significance level} $\alpha$, chosen in advance (often $0.05$):
    if $p \le \alpha$, we \emph{reject} $H_0$ and call the result
    \emph{statistically significant}.
\end{itemize}
\end{definition}

A small p-value means: if $H_0$ were true, data like ours would be surprising. So
either something surprising happened, or $H_0$ is false.

\begin{example}[Is the new model better?]\label{ex:hypothesis-test}
A baseline classifier has accuracy $80\%$. A new model is correct on $340$ of $400$
test examples ($85\%$). Is this convincing evidence that its accuracy is higher
than $80\%$?

\sol Take $H_0$: the new model's accuracy is $p = 0.8$, and $H_1$: $p > 0.8$. Under
$H_0$ the number of correct predictions is binomial with mean $400\times0.8 = 320$
and standard deviation $\sqrt{400\times0.8\times0.2} = 8$. The observed $340$ is
\[
  z = \frac{340 - 320}{8} = 2.5
\]
standard deviations above the mean. By the central limit theorem, the p-value is
approximately $P(Z \ge 2.5) \approx 0.006$ (\cref{fig:p-value}); an exact
binomial computation gives the same. Since $0.006 < 0.05$, we reject $H_0$: the
improvement is very unlikely to be luck.
\end{example}

\begin{figure}[htbp]
\centering
\begin{tikzpicture}
\begin{axis}[mlplot,axis lines=left,width=10cm,height=5cm,xmin=-4,xmax=4,ymin=0,ymax=0.45,
  xtick={-3,...,3},
  ytick=\empty,grid=none,clip=false,xlabel={$z$},
  xlabel style={at={(ticklabel cs:0.5)},anchor=north}]
  \addplot[name path=f,MLdef,very thick,domain=-4:4,samples=150]{exp(-x^2/2)/sqrt(2*pi)};
  \path[name path=axis] (axis cs:-4,0) -- (axis cs:4,0);
  \addplot[PlotRed!50] fill between[of=f and axis,soft clip={domain=2.5:4}];
  \draw[PlotRed,thick] (axis cs:2.5,0) -- (axis cs:2.5,0.13);
  \node[PlotRed,font=\small,anchor=south west] at (axis cs:2.35,0.13) {observed $z = 2.5$};
  \draw[-{Stealth[length=1.6mm]},PlotRed] (axis cs:3.45,0.06) -- (axis cs:2.85,0.008);
  \node[PlotRed,font=\small,anchor=south west] at (axis cs:3.35,0.06) {$p \approx 0.006$};
\end{axis}
\end{tikzpicture}
\caption{The p-value of \cref{ex:hypothesis-test} is the area of the tail of the
  standard normal distribution beyond the observed test statistic.}
\label{fig:p-value}
\end{figure}

If we had asked whether the new model is \emph{different} (better or worse), we
would count extreme results in both directions; this \emph{two-sided} p-value is
$P(\abs Z \ge 2.5) \approx 0.012$.

\subsection{Two kinds of error}

A test can go wrong in two ways:
\begin{center}
\begin{tabular}{lcc}
\toprule
 & $H_0$ true & $H_0$ false\\
\midrule
reject $H_0$ & \textbf{Type I error} (false positive) & correct\\
do not reject & correct & \textbf{Type II error} (false negative)\\
\bottomrule
\end{tabular}
\end{center}
\index{Type I error}\index{Type II error}
The significance level $\alpha$ is the probability of a Type I error. The
probability of correctly rejecting a false $H_0$ is the \term{power} of the test;
it grows with the amount of data. This is the same trade-off as between false
positives and false negatives of a classifier, or of the medical test of
\cref{ex:medical-test}.

\begin{pitfall}[Misreading p-values]
p-values are among the most misinterpreted numbers in science. Keep four points
in mind.
\begin{itemize}
  \item \textbf{A p-value is not the probability that $H_0$ is true.} It is
    $P(\text{data this extreme} \mid H_0)$, not $P(H_0 \mid \text{data})$; confusing
    the two is the error of \cref{sec:conditional}.
  \item \textbf{``Not significant'' does not mean ``no effect''.} With little
    data, a test has low power and may miss a real effect.
  \item \textbf{Significant does not mean important.} With a million test
    examples, an improvement of $0.1$ percentage points can be highly significant
    and practically irrelevant. Report the size of the effect, with a confidence
    interval, not just the p-value.
  \item \textbf{Multiple comparisons.} Test $20$ useless ideas at $\alpha = 0.05$
    and, on average, one will look significant: the probability of at least one
    false positive is $1 - 0.95^{20} \approx 0.64$. Trying many models or
    hyperparameters and reporting the best test score has the same effect.
\end{itemize}
\end{pitfall}

\begin{remark}[Tests and confidence intervals]
A two-sided test at level $5\%$ rejects $H_0\colon\mu = \mu_0$ exactly when (up to
the approximations involved) $\mu_0$ lies outside the $95\%$ confidence interval.
Confidence intervals carry more information, because they also show the size of
the effect and the precision with which it is known.
\end{remark}

\begin{mlnote}[Testing in machine learning]
Statistical tests are used to compare models (preferably \emph{paired} tests that
evaluate both models on the same test examples), to run A/B tests of deployed
systems, and to check whether an apparent improvement exceeds the noise between
random seeds. The multiple-comparisons problem is the reason for keeping a final
test set that is used only once, after all decisions have been made
(\cref{sec:generalisation}).
\end{mlnote}

\sectionexercises

\begin{exercise}
A coin shows $61$ heads in $100$ tosses. Test the hypothesis that it is fair,
using a two-sided test at the $5\%$ level and the normal approximation.
\end{exercise}
\begin{solution}
Under $H_0$ the number of heads has mean $50$ and standard deviation
$\sqrt{100\times0.25} = 5$, so $z = \tfrac{61 - 50}{5} = 2.2$. The two-sided
p-value is $P(\abs Z \ge 2.2) \approx 0.028 < 0.05$: we reject fairness.
\end{solution}

\begin{exercise}
A study reports $p = 0.03$. Which statements are correct? (a) There is a $3\%$
probability that $H_0$ is true. (b) If $H_0$ were true, results at least this
extreme would occur about $3\%$ of the time. (c) The effect is large.
(d) At the $5\%$ level, $H_0$ is rejected.
\end{exercise}
\begin{solution}
Only (b) and (d). Statement (a) confuses $P(\text{data} \mid H_0)$ with
$P(H_0 \mid \text{data})$, and (c) confuses significance with size: a tiny effect
can be significant if there is enough data.
\end{solution}

\begin{exercise}
You compare $20$ hyperparameter settings, none of which is actually better than
the baseline, each with a test at $\alpha = 0.05$ (assume the tests are
independent). What is the probability that at least one comparison looks
significant? What per-test level would keep this probability below about $0.05$?
\end{exercise}
\begin{solution}
$1 - 0.95^{20} \approx 0.64$. Testing each at $\alpha = \tfrac{0.05}{20} = 0.0025$
(the \emph{Bonferroni correction}) gives $1 - 0.9975^{20} \approx 0.049$.
\end{solution}

\begin{exercise}
A model has accuracy $0.9$ on $400$ test examples. Is it significantly different
from a competitor whose accuracy is known to be $0.88$? Use the confidence
interval of the second exercise of \cref{sec:clt}.
\end{exercise}
\begin{solution}
The $95\%$ confidence interval $[0.871, 0.929]$ contains $0.88$, so the difference
is not significant at the $5\%$ level. (Directly: $z = \tfrac{0.9 -
0.88}{\sqrt{0.88\times0.12/400}} \approx 1.23$, with two-sided p-value about $0.22$.)
\end{solution}

%----------------------------------------------------------------------
\section{Generalisation, overfitting and the bias--variance trade-off}\label{sec:generalisation}
%----------------------------------------------------------------------

The goal of learning is not to do well on the training data, which we already
have, but on \emph{new} data from the same distribution: to \term{generalise}.
The training error is an optimistic estimate of the error on new data, because
the parameters were chosen to fit precisely those examples.

\begin{definition}[Training error, test error, generalisation gap]\label{def:generalisation}
The \term{training error} is the average loss of a model on the data it was
trained on. Its \term{test error} (or \emph{generalisation error}) is its expected
loss on new examples from the same distribution, estimated on a held-out
\term{test set}. The difference between the two is the \emph{generalisation gap}.
\end{definition}

\begin{example}[Fitting polynomials]\label{ex:polynomial-overfitting}
Ten points are generated from $y = \sin(2\pi x)$ plus Gaussian noise with
standard deviation $0.25$. Polynomials of degree $d$ are fitted by least squares
(\cref{sec:projections}), and their root mean squared errors (RMSE) are measured on
the training points and on $2000$ fresh test points:
\begin{center}
\begin{tabular}{lcccc}
\toprule
degree $d$ & $0$ & $1$ & $3$ & $9$\\
\midrule
training RMSE & $0.68$ & $0.55$ & $0.35$ & $0.00$\\
test RMSE & $0.75$ & $0.55$ & $0.30$ & $1.05$\\
\bottomrule
\end{tabular}
\end{center}
The degree-$9$ polynomial has $10$ coefficients, enough to pass exactly through
all $10$ points: zero training error. But it achieves this by wild oscillations
between the points (\cref{fig:overfitting}), and its test error is the worst of
all. The cubic captures the true shape and comes close to the noise level $0.25$,
the best any model can do.
\end{example}

\begin{figure}[htbp]
\centering
\begin{tikzpicture}
\begin{groupplot}[group style={group size=3 by 1,horizontal sep=0.45cm,y descriptions at=edge left},
  mlplot,axis lines=box,width=5.3cm,height=4.6cm,xmin=0,xmax=1,ymin=-2,ymax=2,
  xtick={0,0.5,1},ytick={-2,-1,0,1,2},grid=none,title style={font=\small}]
  \nextgroupplot[title={degree $1$: underfitting}]
    \addplot[gray,dashed,domain=0:1,samples=100]{sin(360*x)};
    \addplot[PlotRed] coordinates {(0.00,0.716) (0.01,0.703) (0.02,0.690) (0.03,0.678) (0.04,0.665) (0.05,0.653) (0.06,0.640)
      (0.07,0.627) (0.08,0.615) (0.09,0.602) (0.10,0.590) (0.11,0.577) (0.12,0.564) (0.13,0.552)
      (0.14,0.539) (0.15,0.526) (0.16,0.514) (0.17,0.501) (0.18,0.489) (0.19,0.476) (0.20,0.463)
      (0.21,0.451) (0.22,0.438) (0.23,0.426) (0.24,0.413) (0.25,0.400) (0.26,0.388) (0.27,0.375)
      (0.28,0.363) (0.29,0.350) (0.30,0.337) (0.31,0.325) (0.32,0.312) (0.33,0.299) (0.34,0.287)
      (0.35,0.274) (0.36,0.262) (0.37,0.249) (0.38,0.236) (0.39,0.224) (0.40,0.211) (0.41,0.199)
      (0.42,0.186) (0.43,0.173) (0.44,0.161) (0.45,0.148) (0.46,0.136) (0.47,0.123) (0.48,0.110)
      (0.49,0.098) (0.50,0.085) (0.51,0.072) (0.52,0.060) (0.53,0.047) (0.54,0.035) (0.55,0.022)
      (0.56,0.009) (0.57,-0.003) (0.58,-0.016) (0.59,-0.028) (0.60,-0.041) (0.61,-0.054)
      (0.62,-0.066) (0.63,-0.079) (0.64,-0.091) (0.65,-0.104) (0.66,-0.117) (0.67,-0.129)
      (0.68,-0.142) (0.69,-0.155) (0.70,-0.167) (0.71,-0.180) (0.72,-0.192) (0.73,-0.205)
      (0.74,-0.218) (0.75,-0.230) (0.76,-0.243) (0.77,-0.255) (0.78,-0.268) (0.79,-0.281)
      (0.80,-0.293) (0.81,-0.306) (0.82,-0.318) (0.83,-0.331) (0.84,-0.344) (0.85,-0.356)
      (0.86,-0.369) (0.87,-0.382) (0.88,-0.394) (0.89,-0.407) (0.90,-0.419) (0.91,-0.432)
      (0.92,-0.445) (0.93,-0.457) (0.94,-0.470) (0.95,-0.482) (0.96,-0.495) (0.97,-0.508)
      (0.98,-0.520) (0.99,-0.533) (1.00,-0.545)};
    \addplot[only marks,mark=*,mark size=1.6pt,PlotBlue] coordinates {(0.000,0.457) (0.111,-0.127) (0.222,1.224) (0.333,0.883) (0.444,0.672) (0.556,-0.246)
      (0.667,-0.409) (0.778,-0.977) (0.889,-0.772) (1.000,0.145)};
  \nextgroupplot[title={degree $3$: a good fit}]
    \addplot[gray,dashed,domain=0:1,samples=100]{sin(360*x)};
    \addplot[PlotRed] coordinates {(0.00,0.037) (0.01,0.127) (0.02,0.211) (0.03,0.290) (0.04,0.364) (0.05,0.433) (0.06,0.498)
      (0.07,0.557) (0.08,0.612) (0.09,0.662) (0.10,0.708) (0.11,0.750) (0.12,0.787) (0.13,0.821)
      (0.14,0.850) (0.15,0.875) (0.16,0.897) (0.17,0.915) (0.18,0.930) (0.19,0.940) (0.20,0.948)
      (0.21,0.952) (0.22,0.953) (0.23,0.952) (0.24,0.947) (0.25,0.939) (0.26,0.929) (0.27,0.916)
      (0.28,0.900) (0.29,0.882) (0.30,0.862) (0.31,0.839) (0.32,0.815) (0.33,0.788) (0.34,0.760)
      (0.35,0.730) (0.36,0.698) (0.37,0.664) (0.38,0.629) (0.39,0.593) (0.40,0.555) (0.41,0.517)
      (0.42,0.477) (0.43,0.436) (0.44,0.395) (0.45,0.352) (0.46,0.310) (0.47,0.266) (0.48,0.222)
      (0.49,0.178) (0.50,0.134) (0.51,0.089) (0.52,0.045) (0.53,0.001) (0.54,-0.043)
      (0.55,-0.087) (0.56,-0.130) (0.57,-0.173) (0.58,-0.215) (0.59,-0.257) (0.60,-0.297)
      (0.61,-0.337) (0.62,-0.375) (0.63,-0.413) (0.64,-0.449) (0.65,-0.483) (0.66,-0.517)
      (0.67,-0.548) (0.68,-0.578) (0.69,-0.606) (0.70,-0.632) (0.71,-0.657) (0.72,-0.679)
      (0.73,-0.699) (0.74,-0.716) (0.75,-0.731) (0.76,-0.744) (0.77,-0.754) (0.78,-0.761)
      (0.79,-0.765) (0.80,-0.766) (0.81,-0.765) (0.82,-0.760) (0.83,-0.752) (0.84,-0.740)
      (0.85,-0.725) (0.86,-0.706) (0.87,-0.684) (0.88,-0.658) (0.89,-0.628) (0.90,-0.594)
      (0.91,-0.556) (0.92,-0.513) (0.93,-0.466) (0.94,-0.415) (0.95,-0.360) (0.96,-0.299)
      (0.97,-0.234) (0.98,-0.164) (0.99,-0.089) (1.00,-0.009)};
    \addplot[only marks,mark=*,mark size=1.6pt,PlotBlue] coordinates {(0.000,0.457) (0.111,-0.127) (0.222,1.224) (0.333,0.883) (0.444,0.672) (0.556,-0.246)
      (0.667,-0.409) (0.778,-0.977) (0.889,-0.772) (1.000,0.145)};
  \nextgroupplot[title={degree $9$: overfitting}]
    \addplot[gray,dashed,domain=0:1,samples=100]{sin(360*x)};
    \addplot[PlotRed] coordinates {(0.00,0.457) (0.01,-1.392) (0.02,-2.491) (0.03,-3.028) (0.04,-3.153) (0.05,-2.991)
      (0.06,-2.639) (0.07,-2.176) (0.08,-1.659) (0.09,-1.134) (0.10,-0.632) (0.11,-0.174)
      (0.12,0.226) (0.13,0.561) (0.14,0.831) (0.15,1.037) (0.16,1.183) (0.17,1.277) (0.18,1.326)
      (0.19,1.339) (0.20,1.323) (0.21,1.287) (0.22,1.237) (0.23,1.179) (0.24,1.120) (0.25,1.062)
      (0.26,1.011) (0.27,0.967) (0.28,0.931) (0.29,0.906) (0.30,0.889) (0.31,0.880) (0.32,0.878)
      (0.33,0.882) (0.34,0.888) (0.35,0.895) (0.36,0.900) (0.37,0.902) (0.38,0.899) (0.39,0.889)
      (0.40,0.870) (0.41,0.842) (0.42,0.805) (0.43,0.757) (0.44,0.700) (0.45,0.634) (0.46,0.559)
      (0.47,0.478) (0.48,0.391) (0.49,0.301) (0.50,0.209) (0.51,0.117) (0.52,0.027)
      (0.53,-0.058) (0.54,-0.137) (0.55,-0.209) (0.56,-0.273) (0.57,-0.326) (0.58,-0.370)
      (0.59,-0.403) (0.60,-0.425) (0.61,-0.439) (0.62,-0.444) (0.63,-0.442) (0.64,-0.435)
      (0.65,-0.425) (0.66,-0.415) (0.67,-0.407) (0.68,-0.404) (0.69,-0.408) (0.70,-0.422)
      (0.71,-0.447) (0.72,-0.486) (0.73,-0.539) (0.74,-0.607) (0.75,-0.690) (0.76,-0.785)
      (0.77,-0.890) (0.78,-1.002) (0.79,-1.116) (0.80,-1.226) (0.81,-1.325) (0.82,-1.405)
      (0.83,-1.458) (0.84,-1.475) (0.85,-1.447) (0.86,-1.366) (0.87,-1.223) (0.88,-1.015)
      (0.89,-0.738) (0.90,-0.393) (0.91,0.011) (0.92,0.463) (0.93,0.938) (0.94,1.405)
      (0.95,1.817) (0.96,2.111) (0.97,2.208) (0.98,2.003) (0.99,1.369) (1.00,0.145)};
    \addplot[only marks,mark=*,mark size=1.6pt,PlotBlue] coordinates {(0.000,0.457) (0.111,-0.127) (0.222,1.224) (0.333,0.883) (0.444,0.672) (0.556,-0.246)
      (0.667,-0.409) (0.778,-0.977) (0.889,-0.772) (1.000,0.145)};
\end{groupplot}
\end{tikzpicture}
\caption{Polynomial fits (red) to ten noisy samples (blue) of $\sin(2\pi x)$
  (dashed). The straight line is too simple, the cubic is about right, and the
  degree-$9$ polynomial passes through every point but oscillates wildly
  between them.}
\label{fig:overfitting}
\end{figure}

\begin{definition}[Underfitting and overfitting]\label{def:overfitting}
A model \term[underfitting]{underfits} when it is too simple to capture the
pattern in the data: both its training and its test error are high. It
\term[overfitting]{overfits} when it is so flexible that it fits the noise in
the training data: its training error is low but its test error is high.
\end{definition}

\begin{figure}[htbp]
\centering
\begin{tikzpicture}
\begin{axis}[mlplot,axis lines=left,width=10cm,height=6cm,xmin=-0.3,xmax=9.3,ymin=0,ymax=1.15,
  xtick={0,...,9},xlabel={polynomial degree $d$ (model complexity)},ylabel={RMSE},
  xlabel style={at={(ticklabel cs:0.5)},anchor=north},
  ylabel style={at={(ticklabel cs:0.5)},anchor=south,rotate=90},
  legend style={font=\small,at={(0.03,0.97)},anchor=north west,draw=gray!40},
  legend cell align=left]
  \addplot[PlotBlue,mark=*,mark size=1.6pt] coordinates
    {(0,0.684) (1,0.553) (2,0.552) (3,0.353) (4,0.284) (5,0.230) (6,0.158) (7,0.119) (8,0.116) (9,0.000)};
  \addlegendentry{training error}
  \addplot[PlotRed,mark=square*,mark size=1.6pt] coordinates
    {(0,0.752) (1,0.547) (2,0.548) (3,0.297) (4,0.351) (5,0.347) (6,0.384) (7,0.425) (8,0.445) (9,1.053)};
  \addlegendentry{test error}
  \addplot[gray,dashed,domain=-0.3:9.3,samples=2]{0.25};
  \node[gray,font=\small,anchor=south east] at (axis cs:9.2,0.25) {noise level};
  \node[font=\small,align=center] at (axis cs:1.2,0.12) {under-\\fitting};
  \node[font=\small,align=center] at (axis cs:7,0.75) {overfitting};
\end{axis}
\end{tikzpicture}
\caption{Training and test error of the polynomial fits of
  \cref{ex:polynomial-overfitting} for every degree. The training error always
  decreases with complexity; the test error first decreases, then increases.}
\label{fig:train-test}
\end{figure}

\subsection{The bias--variance trade-off}

Why does the test error first fall and then rise? Imagine repeating the whole
experiment many times, each time with a fresh training set, and look at the
prediction $\hat f(x)$ at one fixed input $x$. Its error has three sources.

\begin{theorem}[Bias--variance decomposition]\label{thm:bias-variance}
Suppose $y = f(x) + \varepsilon$, where the noise $\varepsilon$ has mean $0$ and
variance $\sigma^2$ and is independent of the training set, and let $\hat f$ be
learned from a random training set. Then the expected squared error of the
prediction at $x$ is
\[
  \E\bigl[(y - \hat f(x))^2\bigr]
  = \underbrace{\bigl(f(x) - \E[\hat f(x)]\bigr)^2}_{\text{bias}^2}
  + \underbrace{\Var\bigl(\hat f(x)\bigr)}_{\text{variance}}
  + \underbrace{\sigma^2}_{\text{noise}} .
\]
\end{theorem}

\begin{proof}
Write $y - \hat f(x) = \varepsilon + \bigl(f(x) - \hat f(x)\bigr)$. The noise is
independent of $\hat f(x)$ and has mean $0$, so the cross term has expectation
$0$ and $\E\bigl[(y - \hat f(x))^2\bigr] = \sigma^2 + \E\bigl[(f(x) - \hat
f(x))^2\bigr]$. The last expectation is the mean squared error of $\hat f(x)$ as
an estimator of $f(x)$, which is bias squared plus variance by
\cref{thm:mse-decomposition}.
\end{proof}

\begin{itemize}
  \item \textbf{Bias} is the systematic error of the model class: even averaged
    over many training sets, a straight line cannot follow a sine curve.
  \item \textbf{Variance} measures how much the learned function changes from one
    training set to another. A degree-$9$ polynomial through ten noisy points
    changes completely with every new sample.
  \item \textbf{Noise} is the irreducible error: no model can predict the random
    part of $y$.
\end{itemize}
\index{bias--variance trade-off}
Simple models have high bias and low variance; flexible models have low bias and
high variance. The best test error lies in between (\cref{fig:bias-variance}).

\begin{figure}[htbp]
\centering
\begin{tikzpicture}[scale=0.82]
  \begin{scope}[shift={(0,0)}]
    \fill[MLdef!6] (0,0) circle (1.5);
    \foreach \rad in {1.5,1.0,0.5}{\draw[MLdef!50] (0,0) circle (\rad);}
    \fill[MLdef!40] (0,0) circle (0.12);
    \fill[PlotRed] (0.05,0.12) circle (1.6pt); \fill[PlotRed] (0.05,-0.20) circle (1.6pt); \fill[PlotRed] (0.14,0.07) circle (1.6pt);
    \fill[PlotRed] (-0.08,0.09) circle (1.6pt); \fill[PlotRed] (0.05,0.04) circle (1.6pt); \fill[PlotRed] (0.00,0.08) circle (1.6pt);
    \fill[PlotRed] (-0.11,-0.02) circle (1.6pt); \fill[PlotRed] (-0.07,0.09) circle (1.6pt); \fill[PlotRed] (0.01,-0.04) circle (1.6pt);
    \fill[PlotRed] (-0.12,-0.04) circle (1.6pt);
  \end{scope}
  \begin{scope}[shift={(3.4,0)}]
    \fill[MLdef!6] (0,0) circle (1.5);
    \foreach \rad in {1.5,1.0,0.5}{\draw[MLdef!50] (0,0) circle (\rad);}
    \fill[MLdef!40] (0,0) circle (0.12);
    \fill[PlotRed] (0.00,-0.12) circle (1.6pt); \fill[PlotRed] (0.54,0.42) circle (1.6pt); \fill[PlotRed] (-1.14,-0.79) circle (1.6pt);
    \fill[PlotRed] (-0.07,-0.18) circle (1.6pt); \fill[PlotRed] (0.09,0.09) circle (1.6pt); \fill[PlotRed] (0.89,-0.47) circle (1.6pt);
    \fill[PlotRed] (-0.16,0.86) circle (1.6pt); \fill[PlotRed] (0.27,0.28) circle (1.6pt); \fill[PlotRed] (-0.22,-0.69) circle (1.6pt);
    \fill[PlotRed] (0.07,0.05) circle (1.6pt);
  \end{scope}
  \begin{scope}[shift={(0,-3.4)}]
    \fill[MLdef!6] (0,0) circle (1.5);
    \foreach \rad in {1.5,1.0,0.5}{\draw[MLdef!50] (0,0) circle (\rad);}
    \fill[MLdef!40] (0,0) circle (0.12);
    \fill[PlotRed] (0.42,0.30) circle (1.6pt); \fill[PlotRed] (0.59,0.26) circle (1.6pt); \fill[PlotRed] (0.59,0.41) circle (1.6pt);
    \fill[PlotRed] (0.61,0.32) circle (1.6pt); \fill[PlotRed] (0.69,0.53) circle (1.6pt); \fill[PlotRed] (0.65,0.28) circle (1.6pt);
    \fill[PlotRed] (0.71,0.32) circle (1.6pt); \fill[PlotRed] (0.73,0.24) circle (1.6pt); \fill[PlotRed] (0.74,0.40) circle (1.6pt);
    \fill[PlotRed] (0.41,0.35) circle (1.6pt);
  \end{scope}
  \begin{scope}[shift={(3.4,-3.4)}]
    \fill[MLdef!6] (0,0) circle (1.5);
    \foreach \rad in {1.5,1.0,0.5}{\draw[MLdef!50] (0,0) circle (\rad);}
    \fill[MLdef!40] (0,0) circle (0.12);
    \fill[PlotRed] (0.62,0.51) circle (1.6pt); \fill[PlotRed] (0.19,-0.07) circle (1.6pt); \fill[PlotRed] (0.68,0.20) circle (1.6pt);
    \fill[PlotRed] (0.70,0.72) circle (1.6pt); \fill[PlotRed] (-0.09,0.51) circle (1.6pt); \fill[PlotRed] (1.11,0.28) circle (1.6pt);
    \fill[PlotRed] (0.26,0.72) circle (1.6pt); \fill[PlotRed] (0.71,0.78) circle (1.6pt); \fill[PlotRed] (0.46,-0.22) circle (1.6pt);
    \fill[PlotRed] (0.55,0.21) circle (1.6pt);
  \end{scope}
  \node[font=\small] at (0,1.9) {low variance};
  \node[font=\small] at (3.4,1.9) {high variance};
  \node[font=\small,rotate=90] at (-1.95,0) {low bias};
  \node[font=\small,rotate=90] at (-1.95,-3.4) {high bias};
\end{tikzpicture}
\caption{Bias and variance as a game of darts. Each dot is the prediction of a
  model trained on a different training set; the centre is the truth. Bias is
  the distance from the centre of the cloud to the bull's-eye; variance is the
  spread of the cloud.}
\label{fig:bias-variance}
\end{figure}

\subsection{Fighting overfitting}

\begin{itemize}
  \item \textbf{More data.} Variance shrinks as the training set grows; a
    flexible model that overfits $10$ points may be perfect for $10\,000$.
  \item \textbf{Regularisation.} Penalise large weights ($L^2$ or $L^1$,
    \cref{sec:map}), randomly drop units during training (\emph{dropout}), or
    stop training when the validation error starts to rise (\emph{early
    stopping}).
  \item \textbf{A simpler model}, or fewer features.
  \item \textbf{Data augmentation}: create new training examples by
    transformations that do not change the label (rotating or cropping images).
  \item \textbf{Ensembles}: averaging several models reduces variance.
\end{itemize}
Underfitting calls for the opposite: a bigger model, better features, longer
training or less regularisation.

\subsection{Validation sets and cross-validation}

Choices such as the polynomial degree, the regularisation strength $\lambda$ or
the learning rate are called \emph{hyperparameters}. They cannot be chosen by
training error (which always favours complexity), nor by test error (which would
then no longer be an honest estimate: this is the multiple-comparisons problem of
\cref{sec:hypothesis-testing}). The standard solution is to split the data three
ways:
\begin{itemize}
  \item a \textbf{training set}, used to fit the parameters;
  \item a \term{validation set}, used to choose hyperparameters and compare models;
  \item a \textbf{test set}, used once, at the very end, to estimate the
    performance of the final model.
\end{itemize}

When data are scarce, \term{cross-validation} uses them more efficiently. In
\emph{$k$-fold} cross-validation, the data are split into $k$ equal parts
(\emph{folds}); the model is trained $k$ times, each time on $k - 1$ folds and
validated on the remaining one, and the $k$ validation scores are averaged
(\cref{fig:cross-validation}).

\begin{figure}[htbp]
\centering
\begin{tikzpicture}[x=1.45cm,y=0.62cm]
  \foreach \r in {1,...,5}{
    \node[font=\small,anchor=east] at (0,-\r+0.5) {fold \r};
    \foreach \c in {1,...,5}{
      \ifnum\c=\r
        \fill[MLintuition!70,draw=white,thick] (\c-1,-\r) rectangle (\c,-\r+1);
      \else
        \fill[MLdef!35,draw=white,thick] (\c-1,-\r) rectangle (\c,-\r+1);
      \fi}
    \node[font=\small,anchor=west] at (5.1,-\r+0.5) {$\to$ score $s_{\r}$};}
  \fill[MLdef!35] (0.2,-6.1) rectangle (0.5,-5.7);
  \node[font=\small,anchor=west] at (0.55,-5.9) {training};
  \fill[MLintuition!70] (2.0,-6.1) rectangle (2.3,-5.7);
  \node[font=\small,anchor=west] at (2.35,-5.9) {validation};
  \node[font=\small,anchor=west] at (3.6,-5.9) {CV score $= \frac15\sum_k s_k$};
\end{tikzpicture}
\caption{Five-fold cross-validation. Every example is used for validation exactly
  once and for training four times.}
\label{fig:cross-validation}
\end{figure}

\begin{pitfall}[Data leakage]
Any information from the validation or test data that leaks into training makes
the error estimates optimistic. Common leaks: standardising features with the
mean and standard deviation of the \emph{whole} data set instead of the training
set; duplicates or near-duplicates of training examples in the test set; and
choosing a model because of its test score.
\end{pitfall}

\begin{remark}[Overparameterised models and double descent]
Modern neural networks often have far more parameters than training examples
and can fit even randomly shuffled labels perfectly, yet they generalise well on
real data. Beyond the point where a model can interpolate the training data,
the test error can \emph{decrease again} as the model grows, a phenomenon called
\term{double descent}. Explaining this (through the implicit preferences of
gradient descent, the architecture, and the structure of real data) is an active
area of research. The practical tools of this section, held-out data and
regularisation, remain as important as ever.
\end{remark}

\sectionexercises

\begin{exercise}
Diagnose each situation and suggest a remedy. (a) Training error $1\%$,
validation error $30\%$. (b) Training error $25\%$, validation error $26\%$,
while humans achieve $5\%$.
\end{exercise}
\begin{solution}
(a) Overfitting (high variance): get more data, regularise, augment the data or
use a simpler model. (b) Underfitting (high bias): use a bigger model or better
features, train longer, or reduce regularisation.
\end{solution}

\begin{exercise}
In $5$-fold cross-validation on $1000$ examples, how many times is the model
trained, and on how many examples is each trained and validated?
\end{exercise}
\begin{solution}
It is trained $5$ times, each time on $800$ examples, and validated on the
remaining $200$.
\end{solution}

\begin{exercise}
At some input, a model has squared bias $0.04$ and variance $0.10$, and the noise
variance is $0.25$. What is the expected squared prediction error? Which parts
could a better method reduce?
\end{exercise}
\begin{solution}
$0.04 + 0.10 + 0.25 = 0.39$. Bias and variance can be reduced (by a better model
class, more data, regularisation); the noise $0.25$ cannot.
\end{solution}

\begin{exercise}
You train $100$ models and report the test error of the one with the lowest test
error. Why is this number likely to be too optimistic?
\end{exercise}
\begin{solution}
The test errors contain random fluctuations, and choosing the minimum selects the
model whose fluctuation happened to be most favourable, the multiple-comparisons
effect. An honest estimate requires choosing the model on a validation set and
then evaluating it once on an untouched test set.
\end{solution}

\begin{exercise}
Explain, using the bias--variance decomposition, why collecting more training data
usually helps an overfitting model much more than an underfitting one.
\end{exercise}
\begin{solution}
More data reduces the variance of the learned function (it depends less on the
particular sample), which is the main problem of an overfitting model. The bias of
an underfitting model comes from its limited model class and does not go away
with more data.
\end{solution}

%----------------------------------------------------------------------
\section{Chapter summary}\label{sec:ch10-summary}
%----------------------------------------------------------------------

\begin{keytakeaways}
\begin{itemize}
  \item Data are modelled as i.i.d.\ samples from a population. The mean is
    sensitive to outliers; the median is robust. The sample variance divides by
    $n - 1$ to be unbiased.
  \item An estimator has a bias and a variance; $\MSE = \Bias^2 + \Var$. The
    sample mean has standard error $\sigma/\sqrt n$: four times the data halves
    the error.
  \item Law of large numbers: averages converge to expectations, which justifies
    Monte Carlo estimation (mini-batch gradients, test-set evaluation).
  \item Central limit theorem: $\bar X_n \approx \mathcal{N}(\mu, \sigma^2/n)$;
    approximate $95\%$ confidence interval $\bar x \pm 1.96\,s/\sqrt n$.
  \item Maximum likelihood: choose $\theta$ to maximise
    $\sum_i\ln p(x_i \mid \theta)$; for a coin, $\hat\theta = h/n$; for a normal
    distribution, $\hat\mu = \bar x$.
  \item Gaussian noise $\Rightarrow$ squared error; Bernoulli labels
    $\Rightarrow$ binary cross-entropy; categorical labels $\Rightarrow$
    cross-entropy.
  \item MAP = maximum likelihood + prior penalty; a Gaussian prior gives $L^2$
    regularisation, a Laplace prior $L^1$.
  \item A p-value is $P(\text{data this extreme} \mid H_0)$, not the probability
    that $H_0$ is true. Beware of multiple comparisons.
  \item Test error $=$ bias$^2$ $+$ variance $+$ noise. Underfitting is high
    bias; overfitting is high variance. Use validation sets or cross-validation
    to choose hyperparameters, and touch the test set only once.
\end{itemize}
\end{keytakeaways}

We now have all the ingredients of modern deep learning: functions and their
compositions, vectors and matrices, gradients and optimisation, probability and
likelihood. The final chapter puts them together: softmax and cross-entropy,
information theory, and the backpropagation algorithm that trains neural
networks.
